\subsection{超立方体上的多尺度斯托克斯能量与边界采样证书}\label{subsec:conclusion-multiscale-stokes-energy-sampling}
本小节在欧氏胞腔复形上将广义斯托克斯公式提升为一条多尺度能量恒等式：仅由边界积分构成的量在尺度极限中严格重建 \(\|d\alpha\|_{L^2}\)。
进一步地，在每个胞腔边界上作独立随机采样，可得到带统一置信界与样本复杂度的可审计证书。

\paragraph{离散链复形上的 Stokes--Stein 对偶、Lee--Yang 零测度的 Stokes 反演与代数单值化量子化}
本段给出一条以广义斯托克斯机制为核心的统一定理链：在有限加权图（超立方体或其诱导子图）上得到 Stokes--Stein 完备性与最小能流；在 Lee--Yang 零测度极限口径下，把自由能、零点与大偏差谱组织为 Stokes--Laplacian 反演；当大偏差谱由代数分支承载时，多值性被量子化为整数 Stokes 电荷，并强制导出递归型误差证书结构。

\begin{theorem}[有限加权图上的广义斯托克斯恒等式、Stein 完备性与最小能流]\label{thm:conclusion-stokes-stein-weighted-graph-minflow}
设 \(G=(V,E)\) 为有限连通无向图，并取顶点权
\(\pi:V\to(0,1)\)（\(\sum_{x\in V}\pi(x)=1\)）
与对称导通率 \(c:E\to(0,\infty)\)（记 \(c_{xy}=c_{yx}\) 当 \(\{x,y\}\in E\)）。
令
\[
C^0(G):=\{f:V\to\RR\},
\qquad
C^1_{\mathrm{as}}(G):=\{J:V\times V\to\RR:\ J(x,y)=-J(y,x),\ J(x,y)=0\ \text{若}\ \{x,y\}\notin E\}.
\]
定义离散外微分 \(d:C^0(G)\to C^1_{\mathrm{as}}(G)\) 与加权余微分（散度算子）\(\delta:C^1_{\mathrm{as}}(G)\to C^0(G)\)：
\[
(df)(x,y):=f(y)-f(x),
\qquad
(\delta J)(x):=\frac{1}{\pi(x)}\sum_{y:\{x,y\}\in E}c_{xy}\,J(x,y).
\]
并引入内积
\[
\langle f,g\rangle_\pi:=\sum_{x\in V}\pi(x)f(x)g(x),
\qquad
\langle J,K\rangle_c:=\frac12\sum_{\{x,y\}\in E}c_{xy}\,J(x,y)K(x,y).
\]
则：
\begin{enumerate}
  \item 对任意 \(f\in C^0(G)\) 与 \(J\in C^1_{\mathrm{as}}(G)\)，有精确恒等式
  \[
  \langle f,\delta J\rangle_\pi=-\langle df,J\rangle_c.
  \]
  \item 令零均值子空间
  \[
  C^0_0(G):=\{g\in C^0(G):\langle g,\mathbf 1\rangle_\pi=0\}.
  \]
  则
  \[
  \mathrm{Im}(\delta)=C^0_0(G).
  \]
  特别地，\(\langle \delta J,\mathbf 1\rangle_\pi=0\) 对任意 \(J\) 成立。
  \item 令加权拉普拉斯算子 \(L:=\delta d:C^0(G)\to C^0(G)\)。
  对任意 \(g\in C^0_0(G)\)，存在唯一 \(h\in C^0_0(G)\) 满足
  \[
  Lh=g.
  \]
  并且在全部满足 \(\delta J=g\) 的流中，能量
  \(
  \mathcal E(J):=\langle J,J\rangle_c
  \)
  的唯一极小元为
  \[
  J_\star=dh.
  \]
\end{enumerate}
\end{theorem}
\leanverified{paper\_conclusion\_stokes\_stein\_weighted\_graph\_minflow}
\begin{proof}
对任意 \(f,J\) 展开
\[
\langle f,\delta J\rangle_\pi
=\sum_{x\in V}\pi(x)f(x)\frac{1}{\pi(x)}\sum_{y:\{x,y\}\in E}c_{xy}J(x,y)
=\sum_{\{x,y\}\in E}c_{xy}\bigl(f(x)J(x,y)+f(y)J(y,x)\bigr).
\]
由 \(J(y,x)=-J(x,y)\) 得
\[
\langle f,\delta J\rangle_\pi
=\sum_{\{x,y\}\in E}c_{xy}\bigl(f(x)-f(y)\bigr)J(x,y)
=-\frac12\sum_{\{x,y\}\in E}c_{xy}(df)(x,y)J(x,y)
=-\langle df,J\rangle_c,
\]
得 (1)。
\par
由 (1) 取 \(f=\mathbf 1\) 得 \(\langle \delta J,\mathbf 1\rangle_\pi=-\langle d\mathbf 1,J\rangle_c=0\)，故 \(\mathrm{Im}(\delta)\subseteq C^0_0(G)\)。
反向包含：取任意 \(g\in C^0_0(G)\)。选一棵生成树 \(T\subseteq G\) 并取根 \(r\)。
对树边按父子取向写作 \(e=(p\to q)\)，定义
\[
J(p,q):=\frac{1}{c_{pq}}\sum_{x\in \mathrm{subtree}(q)}\pi(x)g(x),
\qquad
J(q,p):=-J(p,q),
\]
对非树边置零。
则对子树求和作一次递推可得对每个非根顶点 \(q\) 有 \((\delta J)(q)=g(q)\)；
根处由 \(\sum_x\pi(x)g(x)=0\) 得同样成立，从而 \(g\in\mathrm{Im}(\delta)\)。
故 \(\mathrm{Im}(\delta)=C^0_0(G)\)，得 (2)。
\par
对 (3)，在仿射空间 \(\{J:\delta J=g\}\) 上极小化严格凸二次型 \(\mathcal E(J)\)。
取拉格朗日函数
\[
\mathcal L(J,\lambda):=\langle J,J\rangle_c+\langle \lambda,\delta J-g\rangle_\pi,
\qquad \lambda\in C^0(G).
\]
对 \(J\) 的一阶变分为零等价于
\(
2J+\delta^\ast\lambda=0
\)。
由 (1) 可知 \(\delta^\ast=-d\)，故临界点满足 \(J=d\lambda\)。
代回约束得 \(L\lambda=\delta d\lambda=g\)。
将 \(\lambda\) 规一到 \(C^0_0(G)\)（加常数不改变 \(d\lambda\)），得到唯一解 \(h\in C^0_0(G)\)，并由严格凸性知极小解唯一且为 \(J_\star=dh\)。证毕。
\end{proof}

\begin{corollary}[Stokes--Stein 范数的对偶表述与最优证据对象]\label{cor:conclusion-stokes-stein-dual-norm}
沿用定理 \ref{thm:conclusion-stokes-stein-weighted-graph-minflow} 的记号。
对任意 \(g\in C^0_0(G)\) 定义
\[
\|g\|_{\mathrm{St}}:=\inf\{\|J\|_c:\ J\in C^1_{\mathrm{as}}(G),\ \delta J=g\},
\qquad
\|J\|_c:=\sqrt{\langle J,J\rangle_c},
\]
并对任意 \(f\in C^0(G)\) 定义 \(\|df\|_c:=\sqrt{\langle df,df\rangle_c}\)。
则成立严格对偶
\[
\|g\|_{\mathrm{St}}
=\sup\{\langle g,f\rangle_\pi:\ f\in C^0(G),\ \|df\|_c\le 1\}.
\]
并且若 \(h\in C^0_0(G)\) 为 Poisson 方程 \(Lh=g\) 的解，则最优值满足
\[
\|g\|_{\mathrm{St}}=\|dh\|_c,
\]
且极值可由 \(f=h/\|dh\|_c\) 达到。
\leanverified{paper\_conclusion\_stokes\_stein\_dual\_norm}
\end{corollary}
\begin{proof}
对任意满足 \(\delta J=g\) 的 \(J\) 与任意 \(f\)，由定理 \ref{thm:conclusion-stokes-stein-weighted-graph-minflow}(1) 得
\[
\langle g,f\rangle_\pi=\langle \delta J,f\rangle_\pi=-\langle J,df\rangle_c
\le \|J\|_c\,\|df\|_c.
\]
若 \(\|df\|_c\le 1\)，则 \(\langle g,f\rangle_\pi\le \|J\|_c\)；对 \(J\) 取下确界得
\[
\sup_{\|df\|_c\le 1}\langle g,f\rangle_\pi\le \|g\|_{\mathrm{St}}.
\]
反向不等式：取 \(h\in C^0_0(G)\) 满足 \(Lh=g\)，并令 \(J_\star=dh\)。
则由 (1) 取 \(f=h\)、\(J=dh\) 得
\[
\langle g,h\rangle_\pi=\langle Lh,h\rangle_\pi=\langle dh,dh\rangle_c=\|dh\|_c^2.
\]
取 \(f=h/\|dh\|_c\) 得 \(\|df\|_c=1\) 且
\(
\langle g,f\rangle_\pi=\|dh\|_c
\)。
另一方面由定义 \(\|g\|_{\mathrm{St}}\le \|dh\|_c\)。
两式合并即得对偶恒等式与取等构造。证毕。
\end{proof}

\begin{theorem}[粗粒化推前下的 Stokes 结构交换律]\label{thm:conclusion-coarsegraining-stokes-beck-chevalley}
沿用定理 \ref{thm:conclusion-stokes-stein-weighted-graph-minflow} 的设定。
给定任意满射 \(\Phi:V\twoheadrightarrow \bar V\)，定义粗粒化顶点权与导通率为
\[
\bar\pi(\bar x):=\sum_{\Phi(x)=\bar x}\pi(x),
\qquad
\bar c_{\bar x\bar y}:=\sum_{\substack{\{x,y\}\in E\\ \Phi(x)=\bar x,\ \Phi(y)=\bar y}}c_{xy}
\qquad(\bar x,\bar y\in\bar V),
\]
并记粗图为 \(\bar G=(\bar V,\bar E)\)，其中 \(\bar E:=\{\{\bar x,\bar y\}:\bar c_{\bar x\bar y}>0,\ \bar x\neq \bar y\}\)。
对任意 \(J\in C^1_{\mathrm{as}}(G)\)，定义其导通率加权推前流 \(\Phi_!J\in C^1_{\mathrm{as}}(\bar G)\) 为
\[
(\Phi_!J)(\bar x,\bar y):=
\begin{cases}
\bar c_{\bar x\bar y}^{-1}
\displaystyle\sum_{\substack{\{x,y\}\in E\\ \Phi(x)=\bar x,\ \Phi(y)=\bar y}}
c_{xy}\,J(x,y),
&\bar c_{\bar x\bar y}>0,\\[2ex]
0,&\bar c_{\bar x\bar y}=0,
\end{cases}
\]
并对任意 \(\bar f\in C^0(\bar G)\) 定义拉回 \(\Phi^\ast\bar f:=\bar f\circ\Phi\in C^0(G)\)。
则对任意 \(\bar f\in C^0(\bar G)\) 与任意 \(J\in C^1_{\mathrm{as}}(G)\) 成立严格交换恒等式：
\[
\langle d(\Phi^\ast\bar f),J\rangle_c=\langle \bar d\,\bar f,\ \Phi_!J\rangle_{\bar c},
\qquad
\langle \Phi^\ast\bar f,\delta J\rangle_\pi=\langle \bar f,\ \bar\delta(\Phi_!J)\rangle_{\bar\pi},
\]
其中 \((\bar d,\bar\delta,\langle\cdot,\cdot\rangle_{\bar\pi},\langle\cdot,\cdot\rangle_{\bar c})\) 按定理
\ref{thm:conclusion-stokes-stein-weighted-graph-minflow} 对 \((\bar G,\bar\pi,\bar c)\) 的同式定义给出。
因此对任意 \(J\) 都有粗层广义斯托克斯恒等式
\[
\langle \bar f,\bar\delta(\Phi_!J)\rangle_{\bar\pi}=-\langle \bar d\,\bar f,\Phi_!J\rangle_{\bar c}.
\]
\end{theorem}
\leanverified{paper\_conclusion\_coarsegraining\_stokes\_beck\_chevalley}
\begin{proof}
由定义展开并按 \((\Phi(x),\Phi(y))\) 分组，有
\begin{align*}
\langle d(\Phi^\ast\bar f),J\rangle_c
&=\frac12\sum_{\{x,y\}\in E}c_{xy}\bigl(\bar f(\Phi(y))-\bar f(\Phi(x))\bigr)J(x,y)\\
&=\frac12\sum_{\{\bar x,\bar y\}\in\bar E}\bigl(\bar f(\bar y)-\bar f(\bar x)\bigr)
\sum_{\substack{\{x,y\}\in E\\ \Phi(x)=\bar x,\ \Phi(y)=\bar y}}c_{xy}\,J(x,y)\\
&=\frac12\sum_{\{\bar x,\bar y\}\in\bar E}\bar c_{\bar x\bar y}\,(\bar d\bar f)(\bar x,\bar y)\,(\Phi_!J)(\bar x,\bar y)
=\langle \bar d\,\bar f,\Phi_!J\rangle_{\bar c},
\end{align*}
得到第一式。
\par
第二式由两次应用 Stokes 恒等式得到：由定理 \ref{thm:conclusion-stokes-stein-weighted-graph-minflow}(1) 有
\[
\langle \Phi^\ast\bar f,\delta J\rangle_\pi=-\langle d(\Phi^\ast\bar f),J\rangle_c
=-\langle \bar d\,\bar f,\Phi_!J\rangle_{\bar c},
\]
而对粗图 \(\bar G\) 再应用同一定理的 (1)（以 \(\bar f,\Phi_!J\) 代入）得
\(
-\langle \bar d\,\bar f,\Phi_!J\rangle_{\bar c}=\langle \bar f,\bar\delta(\Phi_!J)\rangle_{\bar\pi}
\)，从而得到第二式。
最后一式即为粗图上的 Stokes 恒等式。证毕。
\end{proof}

\begin{proposition}[推前约束下的能量下确界：粗能量作为微能量的精确消去]\label{prop:conclusion-coarsegraining-energy-infimum}
在定理 \ref{thm:conclusion-coarsegraining-stokes-beck-chevalley} 的设定下，定义微观与粗观能量
\[
\mathcal E(J):=\langle J,J\rangle_c=\frac12\sum_{\{x,y\}\in E}c_{xy}\,J(x,y)^2,
\qquad
\bar{\mathcal E}(\bar J):=\langle \bar J,\bar J\rangle_{\bar c}
=\frac12\sum_{\{\bar x,\bar y\}\in\bar E}\bar c_{\bar x\bar y}\,\bar J(\bar x,\bar y)^2.
\]
对任意给定 \(\bar J\in C^1_{\mathrm{as}}(\bar G)\)，令
\(
\mathfrak J(\bar J):=\{J\in C^1_{\mathrm{as}}(G):\ \Phi_!J=\bar J\}
\)。
则成立严格等式
\[
\inf_{J\in\mathfrak J(\bar J)}\mathcal E(J)=\bar{\mathcal E}(\bar J).
\]
并且极小元唯一：对每条粗边 \(\{\bar x,\bar y\}\in\bar E\)，其在全部连接纤维 \(\Phi^{-1}(\bar x)\) 与 \(\Phi^{-1}(\bar y)\) 的微边上取常值
\[
J(x,y)=\bar J(\bar x,\bar y)\qquad
\bigl(\{x,y\}\in E,\ \Phi(x)=\bar x,\ \Phi(y)=\bar y\bigr),
\]
且在同一纤维内的边上取 \(J\equiv 0\)。
\leanverified{paper\_conclusion\_coarsegraining\_energy\_infimum}
\end{proposition}
\begin{proof}
固定一条粗边 \(\{\bar x,\bar y\}\in\bar E\)，记其下所有跨纤维微边集合为
\(
E_{\bar x\bar y}:=\{\{x,y\}\in E:\Phi(x)=\bar x,\ \Phi(y)=\bar y\}
\)。
设 \(j_e:=J(x,y)\)（按 \(\bar x\to\bar y\) 的取向记号）且 \(c_e:=c_{xy}\)。
推前约束 \(\Phi_!J=\bar J\) 等价于
\[
\sum_{e\in E_{\bar x\bar y}}c_e\,j_e=\bar c_{\bar x\bar y}\,\bar J(\bar x,\bar y),
\qquad
\bar c_{\bar x\bar y}=\sum_{e\in E_{\bar x\bar y}}c_e.
\]
由 Cauchy--Schwarz 不等式，
\[
\Bigl(\sum_{e\in E_{\bar x\bar y}}c_e\,j_e\Bigr)^2
\le
\Bigl(\sum_{e\in E_{\bar x\bar y}}c_e\Bigr)\Bigl(\sum_{e\in E_{\bar x\bar y}}c_e\,j_e^2\Bigr)
=\bar c_{\bar x\bar y}\sum_{e\in E_{\bar x\bar y}}c_e\,j_e^2,
\]
从而
\[
\frac12\sum_{e\in E_{\bar x\bar y}}c_e\,j_e^2
\ge
\frac12\,\bar c_{\bar x\bar y}\,\bar J(\bar x,\bar y)^2.
\]
对所有粗边求和，并注意同纤维内边不影响推前约束且能量非负，得
\(
\mathcal E(J)\ge \bar{\mathcal E}(\bar J)
\)。
\par
当且仅当每个 \(\{\bar x,\bar y\}\) 上等号成立时取到下界；而 Cauchy--Schwarz 取等当且仅当
\(
j_e
\)
在 \(E_{\bar x\bar y}\) 上为常数。
由推前约束该常数必为 \(\bar J(\bar x,\bar y)\)。
同时同纤维内边取 \(J\equiv 0\) 是达到全局最小能量的必要条件。
于是给出唯一极小元并实现 \(\inf\mathcal E=\bar{\mathcal E}\)。证毕。
\end{proof}

\begin{remark}[折叠的具体化与粗层导通图]\label{rem:conclusion-fold-stokes-coarsegraining}
在本文的折叠语义中可取 \(\Phi=\Fold_m:\Omega_m=\{0,1\}^m\twoheadrightarrow X_m\)。
若以 \(\Omega_m\) 上的超立方体骨架为微观图 \(G\)，并取自然权 \((\pi,c)\)，
则定理 \ref{thm:conclusion-coarsegraining-stokes-beck-chevalley} 诱导出仅由折叠纤维与跨纤维边计数决定的粗层导通图 \((X_m,\bar c)\)，并在其上得到严格的 Stokes 交换律。
另一方面，附录 \ref{app:op-algebra-fold-subalgebra} 中 \(L^2(\mathcal A_m,\tau_m)\) 的纤维正交分解与 Jones 投影在每条纤维上为秩一平均投影的结构，为上述“粗层算子是微层算子的函子像”提供了与基本构造一致的算子代数实现接口。
\end{remark}

\begin{theorem}[Lee--Yang 零测度诱导的大偏差谱 Stokes 反演]\label{thm:conclusion-leyang-zero-measure-stokes-inversion-ldp}
令 \(\{Z_m(u)\}_{m\ge 1}\) 为非零复多项式序列，\(Z_m\) 的零点（按重数）记为 \(\{\zeta_{m,j}\}_{j=1}^{N_m}\)，并定义归一化零计数测度
\[
\nu_m:=\frac{1}{m}\sum_{j=1}^{N_m}\delta_{\zeta_{m,j}}.
\]
假设存在有限测度 \(\nu\) 与常数 \(C\in\RR\) 使得 \(\nu_m\Rightarrow \nu\)（弱收敛），并且存在开集
\(U\subset\CC\)
（含某条正实射线的邻域），使得对 \(u\in U\) 有极限
\[
F(u):=\lim_{m\to\infty}\frac{1}{m}\log|Z_m(u)|,
\]
且 \(F\) 在 \(U\setminus \mathrm{supp}(\nu)\) 上为 \(C^2\)，并满足对数势表示
\[
F(u)=\int_{\CC}\log|u-\zeta|\,\dd\nu(\zeta)+C
\qquad (u\in U\setminus \mathrm{supp}(\nu)).
\]
令 \(\Lambda(\theta):=F(e^\theta)\)，并假设 \(\Lambda\) 在某实区间上严格凸且可微。
令 Legendre--Fenchel 变换
\(
I(s):=\sup_{\theta}(s\theta-\Lambda(\theta))
\)
在可达斜率区间内可微。
则：
\begin{enumerate}
  \item 对任意允许的实 \(\theta\)（令 \(u=e^\theta\)），有斜率公式
  \[
  \Lambda'(\theta)=\int_{\CC}\Re\!\left(\frac{u}{u-\zeta}\right)\dd\nu(\zeta).
  \]
  \item 对任意内点 \(s\)（即 \(s\in\mathrm{Int}(\Lambda'(\cdot))\)），存在唯一 \(u(s)>0\) 使得
  \(s=\Lambda'(\log u(s))\)，并且
  \[
  I'(s)=\log u(s),
  \qquad
  I(s)=s\log u(s)-F(u(s)).
  \]
  \item 在分布意义下有
  \[
  \frac{1}{2\pi}\Delta F=\nu\qquad\text{于 }U,
  \]
  即对任意 \(\varphi\in C_c^\infty(U)\) 有
  \(
  \int_U \varphi\,\dd\nu=\frac{1}{2\pi}\int_U F\,\Delta\varphi\,\dd A
  \)。
\end{enumerate}
\end{theorem}
\begin{proof}
在 \(u\in U\setminus\mathrm{supp}(\nu)\) 处，对 \(u\) 微分得到
\(
\partial_u\log|u-\zeta|=\Re\frac{1}{u-\zeta}
\)，故
\(
uF'(u)=\int \Re\frac{u}{u-\zeta}\dd\nu
\)，从而 \(\Lambda'(\theta)=uF'(u)\) 给出 (1)。
\par
严格凸性保证 \(\Lambda'\) 在该区间严格单调，故可逆；令 \(u(s)=e^{\theta(s)}\) 为其逆像即可得唯一性与 \(I'(s)=\theta(s)=\log u(s)\)。
由对偶恒等式 \(I(s)=s\theta(s)-\Lambda(\theta(s))\) 得 (2)。
\par
由基本分布恒等式 \(\Delta\log|u-\zeta|=2\pi\delta_\zeta\)，对 \(\nu\) 积分并交换 \(\Delta\) 与积分（以测试函数定义）即得 (3)。证毕。
\end{proof}

\begin{theorem}[代数谱曲线下的大偏差势单值化障碍与 Stokes 量子化]\label{thm:conclusion-algebraic-ldp-monodromy-stokes-quantization}
设 \(u(s)\) 为在 \(\CC\setminus B\) 上解析的代数函数分支（\(B\subset\CC\) 有限），并假设 \(u(s)\neq 0\)。
令 \(\theta(s):=\log u(s)\) 为某个局部对数分支，并设在某连通支上存在解析函数 \(I\) 满足
\[
I'(s)=\theta(s).
\]
则对任意基点 \(s_0\in\CC\setminus B\) 与任意闭路 \(\gamma\subset\CC\setminus B\)（以 \(s_0\) 为起讫点），沿 \(\gamma\) 解析延拓后有：
\begin{enumerate}
  \item 存在整数 \(k(\gamma)\in\ZZ\) 使得
  \[
  \theta_{\mathrm{after}}(s_0)=\theta_{\mathrm{before}}(s_0)+2\pi\mathrm{i}\,k(\gamma).
  \]
  并且若 \(\widetilde\gamma\) 为 \(\gamma\) 在代数曲线的解析分支覆盖上的提升闭路，则
  \[
  k(\gamma)=\frac{1}{2\pi\mathrm{i}}\oint_{\widetilde\gamma}\frac{\dd u}{u}.
  \]
  \item 相应地存在常数 \(C(\gamma)\in\CC\) 使得
  \[
  I_{\mathrm{after}}(s)=I_{\mathrm{before}}(s)+2\pi\mathrm{i}\,k(\gamma)\,s+C(\gamma)
  \]
  在 \(s_0\) 的邻域内成立（作为解析延拓后的分支恒等式）。
  \item 二阶导数满足
\[
  I''(s)=\frac{u'(s)}{u(s)}
  \]
  在 \(\CC\setminus B\) 上为单值解析函数，并且仍为代数函数。
\end{enumerate}
\leanverified{paper\_conclusion\_algebraic\_ldp\_monodromy\_stokes\_quantization}
\end{theorem}
\begin{proof}
由于 \(u(s)\neq 0\)，沿提升闭路 \(\widetilde\gamma\) 的像落在 \(\CC^\times\)。
对数微分的周期
\(
\frac{1}{2\pi\mathrm{i}}\oint_{\widetilde\gamma}\frac{\dd u}{u}
\)
等于该像绕原点的绕数，故为整数，得到 (1)。
\par
由 \(I'(s)=\theta(s)\)，延拓后导数改变量为常数 \(2\pi\mathrm{i}k(\gamma)\)，积分得 (2)。
\par
对数分支加常数不改变导数，故 \(I''=(\log u)'\) 单值；又 \(u\) 代数则 \(u'/u\) 属于代数函数域中的有理函数，从而仍代数，得 (3)。证毕。
\end{proof}

\leanverified{paper\_conclusion\_algebraic\_forces\_precursive\_error\_control}
\begin{proposition}[代数性强制的递归结构与高阶误差证书化]\label{prop:conclusion-algebraic-forces-precursive-error-control}
在定理 \ref{thm:conclusion-algebraic-ldp-monodromy-stokes-quantization} 的设定下，
设 \(u(s)\) 满足不可约代数方程
\[
P(u,s)=0,\qquad P\in\CC[u,s],\qquad \deg_u P=d\ge 2,
\]
且 \(u(s)\neq 0\)。
令
\(
y(s):=I''(s)=\frac{u'(s)}{u(s)}
\)。
则：
\begin{enumerate}
  \item 在代数曲线 \(P(u,s)=0\) 上成立有理恒等式
  \[
  y(s)=-\frac{P_s(u,s)}{u\,P_u(u,s)},
  \]
  从而 \(y(s)\) 为代数函数。
  \item 存在不全为零的多项式系数 \(a_k(s)\in\CC[s]\)（\(0\le k\le d\)）使得 \(y\) 满足线性常微分方程
  \[
  \sum_{k=0}^{d}a_k(s)\,y^{(k)}(s)=0.
  \]
  因此对任一正则点 \(s_\ast\notin B\)，\(y\) 的 Taylor 系数序列为 \(P\)-recursive：满足具多项式系数的有限阶线性递推。
\end{enumerate}
\end{proposition}
\begin{proof}
对恒等式 \(P(u(s),s)=0\) 求导得 \(P_u(u,s)u'(s)+P_s(u,s)=0\)，除以 \(u(s)\,P_u(u,s)\) 得 (1)。
\par
令 \(K:=\CC(s,u)\)。由 \(\deg_u P=d\) 得 \([K:\CC(s)]\le d\)。
由 (1) 知 \(y\in K\)，且对代数扩张内元素求导仍落在 \(K\) 中，故 \(y^{(k)}\in K\) 对所有 \(k\ge 0\) 成立。
于是 \(\{y,y',\dots,y^{(d)}\}\) 作为 \(\CC(s)\)-向量组必线性相关，存在 \(b_k(s)\in\CC(s)\) 不全为零使
\(\sum_{k=0}^{d}b_k(s)y^{(k)}=0\)。
通分清分母即可取 \(a_k(s)\in\CC[s]\) 得所示方程，进而得到 \(P\)-recursive 性。证毕。
\end{proof}

\begin{conjecture}[Stokes--Lee--Yang 量子化电荷的 Langlands 提升]\label{conj:conclusion-stokes-leyang-charge-langlands}
设一族有限状态传递结构的配分函数在热力学极限中由某条代数谱曲线的支配分支决定，使得 Lee--Yang 分支集与坏约化素集合在算术意义下可同时编码。
则猜想：定理 \ref{thm:conclusion-algebraic-ldp-monodromy-stokes-quantization} 中的整数 Stokes 电荷 \(k(\gamma)\) 可提升为某个代数环面或半阿贝尔退化几何所关联的 Langlands 参数的同调分量，
并且围绕分支集的单值化障碍在复解析与 \(\ell\)-进单值化中一致。
\par
该提升的可证伪点在于：若存在两种实现使得 Lee--Yang 分支单值化电荷一致而坏约化处的局部因子不一致，则该提升失败；反之，若在可计算实例族中呈现一致性，则支持该提升的结构必然性。
\end{conjecture}

\endinput


\begin{theorem}[Stokes--Legendre 势的仿射单值化表示]\label{thm:conclusion-algebraic-ldp-affine-monodromy-representation}
\leanverified{paper\_conclusion\_algebraic\_ldp\_affine\_monodromy\_representation}
沿用定理 \ref{thm:conclusion-algebraic-ldp-monodromy-stokes-quantization} 的记号。定义
\[
\operatorname{Aff}_s:=\{as+b:\ a,b\in\CC\}
\]
并将其视为加法群。对任意闭路 \(\gamma\subset \CC\setminus B\)，定义
\[
\rho(\gamma):=2\pi i\,k(\gamma)\,s+C(\gamma)\in \operatorname{Aff}_s.
\]
则
\[
\rho:\pi_1(\CC\setminus B,s_0)\to \operatorname{Aff}_s
\]
是一个群同态；更精确地，
\[
k(\gamma_1\gamma_2)=k(\gamma_1)+k(\gamma_2),
\qquad
C(\gamma_1\gamma_2)=C(\gamma_1)+C(\gamma_2).
\]
因此，\(I\) 的全部多值性并不落在一般解析群中，而是严格压缩到一个平移型仿射表示。
\end{theorem}
\begin{proof}
由定理 \ref{thm:conclusion-algebraic-ldp-monodromy-stokes-quantization}，沿闭路 \(\gamma\) 延拓后有
\[
I_{\mathrm{after}}(s)=I_{\mathrm{before}}(s)+2\pi i\,k(\gamma)\,s+C(\gamma).
\]
先沿 \(\gamma_2\) 再沿 \(\gamma_1\) 延拓，则
\[
I\mapsto I+2\pi i\,k(\gamma_2)s+C(\gamma_2)
\mapsto
I+2\pi i\,(k(\gamma_1)+k(\gamma_2))s+C(\gamma_1)+C(\gamma_2).
\]
另一方面，沿复合闭路 \(\gamma_1\gamma_2\) 延拓给出
\[
I\mapsto I+2\pi i\,k(\gamma_1\gamma_2)s+C(\gamma_1\gamma_2).
\]
比较系数即得
\[
k(\gamma_1\gamma_2)=k(\gamma_1)+k(\gamma_2),
\qquad
C(\gamma_1\gamma_2)=C(\gamma_1)+C(\gamma_2).
\]
故 \(\rho\) 为群同态。证毕。
\end{proof}

\begin{theorem}[大偏差势单值化的充要判据]\label{thm:conclusion-algebraic-ldp-singlevaluedness-criterion}
\leanverified{paper\_conclusion\_algebraic\_ldp\_singlevaluedness\_criterion}
以下命题等价：
\begin{enumerate}
  \item \(I\) 在 \(\CC\setminus B\) 上可取为单值解析函数；
  \item 对一切闭路 \(\gamma\subset \CC\setminus B\)，都有
  \[
  k(\gamma)=0,
  \qquad
  C(\gamma)=0;
  \]
  \item \(\theta=I'\) 单值，且任意延拓常数周期都消失。
\end{enumerate}
\end{theorem}
\begin{proof}
\textup{(1)\(\Rightarrow\)(2)} 若 \(I\) 单值，则任意闭路延拓后
\[
I_{\mathrm{after}}=I_{\mathrm{before}}.
\]
由定理 \ref{thm:conclusion-algebraic-ldp-monodromy-stokes-quantization}，
\[
I_{\mathrm{after}}-I_{\mathrm{before}}=2\pi i\,k(\gamma)\,s+C(\gamma).
\]
右端作为关于 \(s\) 的仿射函数恒等于零，故必有
\[
k(\gamma)=0,
\qquad
C(\gamma)=0.
\]
\textup{(2)\(\Rightarrow\)(1)} 若所有闭路上的延拓增量都为零，则解析延拓与路径无关，故 \(I\) 全局单值。

\textup{(2)\(\Leftrightarrow\)(3)} 由定理 \ref{thm:conclusion-algebraic-ldp-monodromy-stokes-quantization}，
\[
\theta_{\mathrm{after}}-\theta_{\mathrm{before}}=2\pi i\,k(\gamma).
\]
故 \(\theta\) 单值当且仅当 \(k\equiv 0\)。在此条件下，\(I\) 的剩余延拓障碍恰为常数周期 \(C(\gamma)\)。证毕。
\end{proof}

\begin{theorem}[二阶微分是消灭全部多值性的最小阶]\label{thm:conclusion-algebraic-ldp-second-derivative-minimal-singlevaluedness}
\leanverified{paper\_conclusion\_algebraic\_ldp\_second\_derivative\_minimal\_singlevaluedness}
函数
\[
y(s)=I''(s)=\frac{u'(s)}{u(s)}
\]
总是单值解析于 \(\CC\setminus B\)，且仍为代数函数。若存在某条闭路 \(\gamma\) 满足
\[
k(\gamma)\neq 0,
\]
则 \(I'\) 不可能单值，故二阶求导是消灭 \(I\) 全部多值性的最小微分阶。
\end{theorem}
\begin{proof}
由定理 \ref{thm:conclusion-algebraic-ldp-monodromy-stokes-quantization}，
\[
I''(s)=\frac{u'(s)}{u(s)}
\]
在 \(\CC\setminus B\) 上单值且为代数函数。
若某条闭路上 \(k(\gamma)\neq 0\)，则
\[
\theta_{\mathrm{after}}-\theta_{\mathrm{before}}=2\pi i\,k(\gamma)\neq 0,
\]
故 \(I'=\theta\) 不是单值函数。另一方面，二阶导数已把全部分支差异消去，因此二阶微分是最小消多值阶。证毕。
\end{proof}

\begin{theorem}[大偏差势模 affine 的全局单值化]\label{thm:conclusion-algebraic-ldp-affine-quotient-singlevaluedness}
\leanverified{paper\_conclusion\_algebraic\_ldp\_affine\_quotient\_singlevaluedness}
在商空间
\[
\mathcal O(\CC\setminus B)/\operatorname{Aff}_s
\]
中，\(I\) 的分支类 \([I]\) 是全局单值的。并且该商类完全由代数函数 \(y=I''\) 决定。
\end{theorem}
\begin{proof}
由定理 \ref{thm:conclusion-algebraic-ldp-monodromy-stokes-quantization}，任意两条解析延拓分支之差总是
\[
2\pi i\,k(\gamma)\,s+C(\gamma)\in \operatorname{Aff}_s.
\]
故模去 \(\operatorname{Aff}_s\) 后，所有分支都对应同一商类，因而 \([I]\) 全局单值。

另一方面，若 \(I_1''=I_2''\)，则
\[
(I_1-I_2)''=0,
\]
故 \(I_1-I_2\in \operatorname{Aff}_s\)。因此商类 \([I]\) 与 \(I''=y\) 一一对应。证毕。
\end{proof}

\begin{theorem}[高阶 Stokes 纠正项的 \texorpdfstring{$d$}{d}-维 D-finite 压缩]\label{thm:conclusion-algebraic-ldp-dfinite-stokes-compression}
\leanverified{paper\_conclusion\_algebraic\_ldp\_dfinite\_stokes\_compression}
设 \(u(s)\) 满足不可约代数方程
\[
P(u,s)=0,
\qquad
\deg_u P=d\ge 2.
\]
则 \(y=I''=u'/u\) 满足一个阶数至多为 \(d\) 的线性常微分方程
\[
\sum_{k=0}^{d}a_k(s)\,y^{(k)}(s)=0,
\qquad
a_k(s)\in \CC[s].
\]
并因此对每个 \(n\ge d\)，存在有理函数 \(r_{n,0},\dots,r_{n,d-1}\in \CC(s)\) 使
\[
y^{(n)}(s)=\sum_{j=0}^{d-1}r_{n,j}(s)\,y^{(j)}(s).
\]
因此，全部高阶 Stokes 纠正项都被压缩进一个至多 \(d\)-维的 D-finite 模中。
\end{theorem}
\begin{proof}
由命题 \ref{prop:conclusion-algebraic-forces-precursive-error-control}，\(y\) 满足阶数至多为 \(d\) 的线性常微分方程。设其最小阶为 \(r\le d\)。则对任意 \(n\ge r\)，将该方程反复求导并消元，即可把 \(y^{(n)}\) 写成
\[
y,\ y',\ \dots,\ y^{(r-1)}
\]
的 \(\CC(s)\)-线性组合。由于 \(r\le d\)，补足到 \(d\)-维生成系不影响结论。证毕。
\end{proof}

\begin{corollary}[局部 \texorpdfstring{$d$}{d}-jet 对商类 \texorpdfstring{$[I]$}{[I]} 的完备决定性]\label{cor:conclusion-algebraic-ldp-local-djet-determines-affine-class}
取任意正则点 \(s_\ast\in \CC\setminus B\)。则商类 \([I]\) 的局部解析胚由有限数据
\[
\bigl(y(s_\ast),y'(s_\ast),\dots,y^{(d-1)}(s_\ast)\bigr)
\]
唯一决定。等价地，\([I]\) 的全部高阶导数在 \(s_\ast\) 处都由上述 \(d\)-jet 递归恢复。
\end{corollary}
\begin{proof}
由定理 \ref{thm:conclusion-algebraic-ldp-dfinite-stokes-compression}，\(y\) 满足阶数至多为 \(d\) 的线性 ODE。经典 ODE 唯一性定理表明，给定
\[
y(s_\ast),\dots,y^{(d-1)}(s_\ast)
\]
即可唯一确定 \(y\) 的局部解，从而唯一确定 \(I''\) 的局部胚。
再由定理 \ref{thm:conclusion-algebraic-ldp-affine-quotient-singlevaluedness}，\([I]\) 与 \(I''\) 一一对应，故结论成立。证毕。
\end{proof}

\begin{theorem}[Stokes 斜率量子群的最小量子]\label{thm:conclusion-algebraic-ldp-stokes-slope-quantization}
\leanverified{paper\_conclusion\_algebraic\_ldp\_stokes\_slope\_quantization}
映射
\[
k:\pi_1(\CC\setminus B,s_0)\to \ZZ
\]
的像是一个循环子群，故存在唯一整数
\[
m_\ast\in \ZZ_{\ge 0}
\]
使
\[
\operatorname{Im}(k)=m_\ast \ZZ.
\]
于是所有分支斜率跳跃都被量子化为
\[
2\pi i\,m_\ast\ZZ.
\]
其中
\[
m_\ast=0
\iff
I' \text{ 单值}.
\]
\end{theorem}
\begin{proof}
由定理 \ref{thm:conclusion-algebraic-ldp-affine-monodromy-representation}，\(k\) 是群同态。由于 \(\ZZ\) 的任意子群都形如 \(m_\ast \ZZ\)，故存在唯一 \(m_\ast\ge 0\) 满足
\[
\operatorname{Im}(k)=m_\ast \ZZ.
\]
于是
\[
\theta_{\mathrm{after}}-\theta_{\mathrm{before}}=2\pi i\,k(\gamma)
\]
总落在 \(2\pi i\,m_\ast\ZZ\) 中。
最后，
\[
m_\ast=0
\iff
\operatorname{Im}(k)=\{0\}
\iff
k\equiv 0,
\]
再由定理 \ref{thm:conclusion-algebraic-ldp-singlevaluedness-criterion}，这与 \(I'\) 单值等价。证毕。
\end{proof}

\begin{corollary}[去对数单值化的最小 covering 与常数余单值]\label{cor:conclusion-algebraic-ldp-minimal-covering-kernel-k}
设
\[
\widetilde X\to \CC\setminus B
\]
是对应于子群 \(\ker k\subset \pi_1(\CC\setminus B,s_0)\) 的 covering。则：
\begin{enumerate}
  \item \(I'\) 在 \(\widetilde X\) 上单值；
  \item \(\widetilde X\) 在“使 \(I'\) 单值”的覆盖中是最小的；
  \item 在 \(\widetilde X\) 上，\(I\) 的剩余多值性只表现为常数平移
  \[
  I_{\mathrm{after}}-I_{\mathrm{before}}=C(\gamma),
  \qquad
  \gamma\in \ker k.
  \]
\end{enumerate}
\end{corollary}
\begin{proof}
由 covering--monodromy 对应，\(\widetilde X\) 上单值对象恰对应那些在 \(\ker k\) 上 monodromy 平凡的局部函数。对 \(I'\) 而言，其 monodromy 为
\[
I'_{\mathrm{after}}-I'_{\mathrm{before}}=2\pi i\,k(\gamma).
\]
因此当且仅当 \(\gamma\in \ker k\) 时该差为零，故 \(I'\) 在 \(\widetilde X\) 上单值。

若某个更小的覆盖也使 \(I'\) 单值，则其对应子群必包含在 \(\ker k\) 内；由覆盖对应的反向序关系可知 \(\widetilde X\) 是最小者。

最后，在 \(\ker k\) 上，定理 \ref{thm:conclusion-algebraic-ldp-monodromy-stokes-quantization} 的仿射跳跃简化为
\[
I_{\mathrm{after}}-I_{\mathrm{before}}=C(\gamma),
\]
即仅剩常数余单值。证毕。
\end{proof}


\begin{theorem}[多尺度斯托克斯能量的 \(L^2\) 重建]\label{thm:conclusion-multiscale-stokes-energy-l2-reconstruction}
令 \(Q=[0,1]^d\)（\(d\ge 2\)），取 \(\alpha\in C^{1}(Q;\Lambda^{k})\)（\(0\le k\le d-1\)）。
对 dyadic 级别 \(m\in\NN\)，记网格步长 \(h:=2^{-m}\)，并令 \(\mathcal C_m^{k+1}\) 为由该网格诱导的全体定向 \((k+1)\)-维小立方体胞腔集合。
对 \(C\in\mathcal C_m^{k+1}\)，定义尺度 \(m\) 的斯托克斯差分
\[
\mathbf S_m\alpha(C)
:=\frac{1}{h^{k+1}}\int_{\partial C}\alpha
=\frac{1}{h^{k+1}}\int_C d\alpha,
\]
其中第二个等号为胞腔上的广义斯托克斯定理。
并定义多尺度斯托克斯能量
\[
\mathcal E_m(\alpha):=h^d\sum_{C\in\mathcal C_m^{k+1}}\bigl|\mathbf S_m\alpha(C)\bigr|^2.
\]
则：
\begin{enumerate}
  \item 若 \(d\alpha\) 连续，则
  \[
  \lim_{m\to\infty}\mathcal E_m(\alpha)=\int_Q |d\alpha(x)|^2\,\mathrm dx.
  \]
  \item 若 \(d\alpha\) 为 Lipschitz，设 \(L:=\mathrm{Lip}(d\alpha)\)，则存在仅依赖 \((d,k)\) 的常数 \(C_{d,k}<\infty\) 使
  \[
  \biggl|\mathcal E_m(\alpha)-\int_Q |d\alpha(x)|^2\,\mathrm dx\biggr|
  \le C_{d,k}\,L\,h.
  \]
\end{enumerate}
\end{theorem}

\begin{proof}
对每个胞腔 \(C\in\mathcal C_m^{k+1}\) 令其 \((k{+}1)\)-体积为 \(|C|=h^{k+1}\)，并记
\[
\overline{d\alpha}_C:=\frac1{|C|}\int_C d\alpha.
\]
由定义与斯托克斯公式有 \(\mathbf S_m\alpha(C)=\overline{d\alpha}_C\)，故
\[
\mathcal E_m(\alpha)
=h^d\sum_{C\in\mathcal C_m^{k+1}}\bigl|\overline{d\alpha}_C\bigr|^2.
\]
将 \(d\alpha\) 在每个 \(C\) 上用其平均 \(\overline{d\alpha}_C\) 近似，可将右端视为对函数 \(|d\alpha|^2\) 的 Riemann 型求和。
当 \(d\alpha\) 连续时，胞腔平均在 \(h\to 0\) 时对几乎处处点一致逼近 \(d\alpha\)；
结合一致连续性与主导收敛，即得 (1)。

若 \(d\alpha\) 为 Lipschitz，则对任意 \(x\in C\) 有
\(
|d\alpha(x)-\overline{d\alpha}_C|\le L\,\mathrm{diam}(C)\le L\sqrt d\,h
\)。
用恒等式
\(
\bigl||u|^2-|v|^2\bigr|\le (|u|+|v|)\,|u-v|
\)
并在 \(C\) 上积分、再对所有 \(C\) 求和，可得整体误差被 \(O(Lh)\) 控制；
把维数与组合常数吸收到 \(C_{d,k}\) 中即得 (2)。证毕。
\end{proof}

\paragraph{Hilbert 投影结构与尺度无关的采样集中}
多尺度斯托克斯能量不仅给出 \(\|d\alpha\|_{L^2}\) 的尺度极限，还携带一个内生的 \(L^2\) 投影/鞅结构，从而导出尺度单调性与正交增量分解；在有界性假设下，胞腔层的 Monte Carlo 估计具备与尺度无关的统一置信界。

\begin{theorem}[多尺度斯托克斯能量的 \(L^2\) 投影结构与正交增量]\label{thm:conclusion-multiscale-stokes-energy-hilbert-projection}
\leanverified{paper\_conclusion\_multiscale\_stokes\_energy\_hilbert\_projection}
沿用定理 \ref{thm:conclusion-multiscale-stokes-energy-l2-reconstruction} 的记号，并令
\(
\beta:=d\alpha
\)。
对固定的终端尺度 \(M\in\NN\)，在有限集合 \(\mathcal C_M^{k+1}\) 上取均匀随机胞腔 \(C_M\)。
对每个 \(0\le m\le M\)，令 \(C_m\in\mathcal C_m^{k+1}\) 为 \(C_M\) 的唯一祖先胞腔（同向并满足 \(C_M\subset C_m\)），并令 \(\mathcal F_m:=\sigma(C_m)\)。
记
\[
X_M:=\mathbf S_M\alpha(C_M)=\frac{1}{h_M^{k+1}}\int_{C_M}\beta,
\qquad h_m:=2^{-m}.
\]
则对每个 \(m\le M\) 有
\begin{equation}\label{eq:conclusion-stokes-energy-conditional-expectation}
\boxed{\;
\EE[X_M\mid \mathcal F_m]=\mathbf S_m\alpha(C_m)\;}
\end{equation}
从而在 \(L^2\) 意义下成立尺度单调性与正交增量分解：
\begin{equation}\label{eq:conclusion-stokes-energy-pythagoras-increments}
\boxed{\;
\EE\bigl[| \mathbf S_{m+1}\alpha(C_{m+1})|^2\bigr]
-\EE\bigl[|\mathbf S_{m}\alpha(C_{m})|^2\bigr]
=\EE\bigl[|\mathbf S_{m+1}\alpha(C_{m+1})-\mathbf S_{m}\alpha(C_{m})|^2\bigr]\ \ge 0.\;}
\end{equation}
\end{theorem}
\begin{proof}
由祖先定义，条件于事件 \(\{C_m=C\}\)（\(C\in\mathcal C_m^{k+1}\)）时，随机变量 \(C_M\) 在 \(C\) 的全部 \(M\)-级子胞腔上均匀分布。
由积分的可加性与 \(\beta\) 在子胞腔上的平均迭代，
\[
\EE[X_M\mid C_m=C]
=\frac{1}{h_M^{k+1}}\cdot \frac{1}{\#\{\text{子胞腔}\}}\sum_{C'\subset C}\int_{C'}\beta
=\frac{1}{h_m^{k+1}}\int_{C}\beta
=\mathbf S_m\alpha(C),
\]
得到 \eqref{eq:conclusion-stokes-energy-conditional-expectation}。
\par
由于 \(\{\mathcal F_m\}_{m\le M}\) 为递增过滤，\(X_m:=\EE[X_M\mid\mathcal F_m]\) 构成 \(L^2\)-鞅。
因此
\(
X_{m+1}-X_m\perp X_m
\)
且 \(\EE[|X_{m+1}|^2]=\EE[|X_m|^2]+\EE[|X_{m+1}-X_m|^2]\)。
把 \(X_m\) 代回 \eqref{eq:conclusion-stokes-energy-conditional-expectation} 即得 \eqref{eq:conclusion-stokes-energy-pythagoras-increments}。证毕。
\end{proof}

\begin{theorem}[顶维 dyadic 分割下的条件期望表征与跨尺度能量恒等式]\label{thm:conclusion-stokes-energy-dyadic-martingale}
\leanverified{paper\_conclusion\_stokes\_energy\_dyadic\_martingale}
沿用定理 \ref{thm:conclusion-multiscale-stokes-energy-l2-reconstruction}，并取顶维情形 \(k=d-1\)。
令 \(\beta=d\alpha\) 并以标准体积形式识别为 \(d\)-维密度 \(f\in L^2(Q)\)：\(\beta=f(x)\,\mathrm dx\)。
令 \(\mathscr F_m\) 为由 dyadic 顶维立方剖分 \(\mathcal C_m^{d}\) 生成的 \(\sigma\)-代数。
则成立：
\begin{enumerate}
  \item 条件期望表征
  \[
  \mathcal E_m(\alpha)=\bigl\|\EE[f\mid \mathscr F_m]\bigr\|_{L^2(Q)}^2.
  \]
  \item 精确方差分解
  \[
  \|f\|_{L^2(Q)}^2-\mathcal E_m(\alpha)
  =\sum_{C\in\mathcal C_m^{d}}\int_C\bigl|f-\langle f\rangle_C\bigr|^2,
  \qquad
  \langle f\rangle_C:=\frac1{|C|}\int_C f=\mathbf S_m\alpha(C),
  \]
  从而 \(m\mapsto \mathcal E_m(\alpha)\) 单调不减且 \(\mathcal E_m(\alpha)\uparrow \|f\|_2^2\)。
  \item 跨尺度能量恒等式与局部 Haar 增量公式
  \[
  \|f\|_2^2=\mathcal E_0(\alpha)+\sum_{m\ge 0}\bigl(\mathcal E_{m+1}(\alpha)-\mathcal E_m(\alpha)\bigr),
  \qquad
  \mathcal E_{m+1}(\alpha)-\mathcal E_m(\alpha)
  =\bigl\|\EE[f\mid\mathscr F_{m+1}]-\EE[f\mid\mathscr F_m]\bigr\|_2^2\ge 0,
  \]
  并且每一增量满足完全显式的局部方差公式
  \[
  \mathcal E_{m+1}(\alpha)-\mathcal E_m(\alpha)
  =\sum_{P\in\mathcal C_m^{d}}\frac{|P|}{2^d}\sum_{\substack{C\subset P\\ C\in\mathcal C_{m+1}^{d}}}
  \bigl|\mathbf S_{m+1}\alpha(C)-\mathbf S_m\alpha(P)\bigr|^2.
  \]
  因此尾项误差具有严格证书表达：
  \(
  \|f\|_2^2-\mathcal E_M(\alpha)=\sum_{m\ge M}(\mathcal E_{m+1}-\mathcal E_m)
  \)。
\end{enumerate}
\end{theorem}
\begin{proof}
对每个 \(C\in\mathcal C_m^{d}\) 记 \(f_C:=\langle f\rangle_C=\frac1{|C|}\int_C f\)。
由顶维斯托克斯读出 \(\mathbf S_m\alpha(C)=h^{-d}\int_C\beta=f_C\)，并且
\(
\EE[f\mid\mathscr F_m]=\sum_{C\in\mathcal C_m^{d}} f_C\,\ind_C
\)。
因此
\[
\bigl\|\EE[f\mid\mathscr F_m]\bigr\|_2^2
=\sum_{C}\int_C |f_C|^2
=\sum_C |C|\,|f_C|^2
=h^d\sum_C |\mathbf S_m\alpha(C)|^2
=\mathcal E_m(\alpha),
\]
得 (1)。
\par
在每个 \(C\) 上有 \(\int_C f_C(f-f_C)=f_C\int_C(f-f_C)=0\)，故
\(
\int_C|f|^2=\int_C|f_C|^2+\int_C|f-f_C|^2
\)。
对 \(C\) 求和并识别第一项为 \(\mathcal E_m(\alpha)\) 即得 (2)；
单调性与 \(\mathcal E_m(\alpha)\uparrow\|f\|_2^2\) 亦可由条件期望范数随细化单调增加推出。
\par
令 \(X_m:=\EE[f\mid\mathscr F_m]\)，则 \((X_m)_{m\ge 0}\) 为 \(L^2\) 鞅且
\(
\|X_{m+1}\|_2^2-\|X_m\|_2^2=\|X_{m+1}-X_m\|_2^2
\)。
代入 (1) 得跨尺度恒等式与非负性。
局部方差公式由逐父胞腔分解：
对每个 \(P\in\mathcal C_m^d\) 记其 \(2^d\) 个子胞腔均值为 \(f_C\) 且 \(f_P\) 为父均值，则
\[
\sum_{C\subset P}|C|\,|f_C|^2-|P|\,|f_P|^2
=\sum_{C\subset P}|C|\,|f_C-f_P|^2,
\]
其中使用 \(\sum_{C\subset P}|C|(f_C-f_P)=0\)。
将 \(|C|=|P|/2^d\) 并代回 \(f_C=\mathbf S_{m+1}\alpha(C)\)、\(f_P=\mathbf S_m\alpha(P)\) 即得所示公式。证毕。
\end{proof}

\begin{corollary}[\(H^1\) 正则性下的二阶尺度误差控制]\label{cor:conclusion-stokes-energy-h1-second-order}
在定理 \ref{thm:conclusion-stokes-energy-dyadic-martingale} 的设定下，若 \(f\in H^1(Q)\)，则存在仅依赖于维数 \(d\) 的常数 \(C_d>0\) 使得对任意 \(m\ge 0\) 有
\[
0\le \|f\|_2^2-\mathcal E_m(\alpha)\le C_d\,2^{-2m}\,\|\nabla f\|_2^2.
\]
\end{corollary}
\begin{proof}
由定理 \ref{thm:conclusion-stokes-energy-dyadic-martingale}(2)，误差为 \(\sum_C\int_C|f-\langle f\rangle_C|^2\)。
对每个边长 \(h=2^{-m}\) 的立方体 \(C\) 用 Poincar\'e 不等式得
\(
\int_C|f-\langle f\rangle_C|^2\le C_d\,h^2\int_C|\nabla f|^2
\)；
求和即得结论。证毕。
\leanverified{paper\_conclusion\_stokes\_energy\_h1\_second\_order}
\end{proof}

\begin{corollary}[顶维胞腔抽样对 \texorpdfstring{$\|d\alpha\|_{L^2}$}{||dα||\_2} 的置信区间]\label{cor:conclusion-stokes-energy-topcell-confidence}
\leanverified{paper\_conclusion\_stokes\_energy\_topcell\_confidence}
在命题 \ref{prop:conclusion-stokes-energy-cell-sampling-concentration} 的设定下取 \(k=d-1\)，并记 \(\widehat E_m:=\widehat{\mathcal E}_m(\alpha)\)。
则对任意 \(0<\delta<1\)，以概率至少 \(1-\delta\) 有
\[
\bigl|\widehat E_m-\|d\alpha\|_{L^2(Q)}^2\bigr|
\le
\bigl(\|d\alpha\|_{L^2(Q)}^2-\mathcal E_m(\alpha)\bigr)
 + M^2\sqrt{\frac{\log(2/\delta)}{2N}}.
\]
若进一步 \(f\in H^1(Q)\)，则可用推论 \ref{cor:conclusion-stokes-energy-h1-second-order} 将离散化偏差显式上界为 \(C_d\,2^{-2m}\|\nabla f\|_2^2\)。
\end{corollary}

\begin{proposition}[胞腔抽样下斯托克斯能量的尺度无关集中界]\label{prop:conclusion-stokes-energy-cell-sampling-concentration}
\leanverified{paper\_conclusion\_stokes\_energy\_cell\_sampling\_concentration}
沿用定理 \ref{thm:conclusion-multiscale-stokes-energy-l2-reconstruction} 的记号，并额外假设 \(d\alpha\in L^\infty(Q)\) 且 \(\|d\alpha\|_\infty\le M\)。
固定尺度 \(m\)，从 \(\mathcal C_m^{k+1}\) 中独立等概率抽取 \(N\) 个胞腔 \(C_1,\dots,C_N\)，并定义估计量
\[
\widehat{\mathcal E}_m(\alpha)
:=
h_m^{d}\,|\mathcal C_m^{k+1}|\cdot \frac1N\sum_{j=1}^{N}\bigl|\mathbf S_m\alpha(C_j)\bigr|^2.
\]
则对任意 \(0<\delta<1\)，以概率至少 \(1-\delta\) 有
\[
\boxed{\;
\bigl|\widehat{\mathcal E}_m(\alpha)-\mathcal E_m(\alpha)\bigr|
\le
h_m^{d}\,|\mathcal C_m^{k+1}|\,M^2\sqrt{\frac{\log(2/\delta)}{2N}}\;}
\]
并且在标准 dyadic 约定下（以半开胞腔计数），有恒等式
\(
h_m^{d}\,|\mathcal C_m^{k+1}|=\binom{d}{k+1}
\)，
从而上界与尺度 \(m\) 无关。
\end{proposition}
\begin{proof}
由 Jensen 不等式与 \(\|d\alpha\|_\infty\le M\) 得
\(
|\mathbf S_m\alpha(C)|\le M
\)，从而随机变量
\(
X_j:=|\mathbf S_m\alpha(C_j)|^2
\)
满足 \(X_j\in[0,M^2]\) 且 \(\EE[X_j]=\frac1{|\mathcal C_m^{k+1}|}\sum_{C}|\mathbf S_m\alpha(C)|^2\)。
由 Hoeffding 不等式，
\[
\PP\!\left(\left|\frac1N\sum_{j=1}^N X_j-\EE[X_1]\right|>\varepsilon\right)
\le 2\exp\!\left(-\frac{2N\varepsilon^2}{M^4}\right).
\]
取 \(\varepsilon=M^2\sqrt{\log(2/\delta)/(2N)}\) 并乘以 \(h_m^{d}|\mathcal C_m^{k+1}|\) 即得结论；计数恒等式由 dyadic 网格的平移不变计数直接给出。证毕。
\end{proof}

\begin{conjecture}[折叠动力学诱导的斯托克斯胞腔平均的统一次高斯性]\label{conj:conclusion-stokes-energy-uniform-subgaussian}
在折叠语法与纤维重构图呈 CAT(0) 立方几何的体系中，取一族由可见相位/审计势自然生成的 \(1\)-形式 \(\alpha\)（使 \(d\alpha\in L^2(Q)\)）。
令 \(C\) 为在尺度 \(m\) 上按平移不变方式抽取的随机 \((k+1)\)-胞腔，并记随机变量
\[
Z_m:=\mathbf S_m\alpha(C)-\EE[\mathbf S_m\alpha(C)].
\]
则存在仅依赖于 \((d,k)\) 与局部折叠约束的常数 \(K<\infty\)，使得对所有 \(m\) 与所有 \(t\ge 0\) 都有一致的次高斯尾界
\[
\PP(|Z_m|\ge t)\le 2\exp\!\left(-\frac{t^2}{K\,\|d\alpha\|_{L^2(Q)}^2}\right).
\]
若该猜想成立，则命题 \ref{prop:conclusion-stokes-energy-cell-sampling-concentration} 的尺度无关集中界可在不依赖 \(L^\infty\) 假设下推广为普适现象，并为尺度鲁棒的泛化界与自适应停止准则提供统一概率桥梁。
\end{conjecture}

\endinput


\paragraph{一般次数 dyadic 级联与正则度读出}
将顶维情形的条件期望机制按系数分量作用到 \(d\alpha\) 上，可把多尺度 Stokes 能量的正交增量推广到任意次数，并由此得到完全由尺度增量决定的后验误差证书与 Besov 正则度读出。
\leanverified{paper\_conclusion\_stokes\_energy\_general\_degree\_cascade}

\begin{theorem}[一般次数多尺度 Stokes 能量的正交级联]\label{thm:conclusion-stokes-energy-general-degree-cascade}
沿用定理 \ref{thm:conclusion-multiscale-stokes-energy-l2-reconstruction} 的设定，并把
\[
d\alpha=\sum_{|I|=k+1} f_I(x)\,\mathrm d x_I,
\qquad
f_I\in L^2(Q).
\]
记 \(\mathscr F_m\) 为边长 \(h_m:=2^{-m}\) 的 dyadic 顶维立方剖分 \(\mathcal Q_m\) 所生成的 \(\sigma\)-代数，\(P_m:=\mathbb E[\cdot\mid \mathscr F_m]\) 为按系数逐分量作用的条件期望。则有
\[
\mathcal E_m(\alpha)
=
\sum_{|I|=k+1}\|P_m f_I\|_{L^2(Q)}^2
=
\|P_m(d\alpha)\|_{L^2(Q)}^2.
\]
并且成立精确增量公式
\[
\mathcal E_{m+1}(\alpha)-\mathcal E_m(\alpha)
=
\|P_{m+1}(d\alpha)-P_m(d\alpha)\|_{L^2(Q)}^2
\]
\[
=
\frac{h_m^d}{2^d}\sum_{Q\in\mathcal Q_m}\sum_{Q'\prec Q}
\bigl|
P_{m+1}(d\alpha)|_{Q'}-P_m(d\alpha)|_{Q}
\bigr|^2,
\]
其中 \(Q'\prec Q\) 遍历 \(Q\) 的 \(2^d\) 个 dyadic 子立方体，而 \(P_j(d\alpha)|_Q\) 表示 \(P_j(d\alpha)\) 在 \(Q\) 上的常值 \((k+1)\)-形式。特别地，\(\mathcal E_m(\alpha)\) 随 \(m\) 单调不减。
\end{theorem}
\begin{proof}
由定理 \ref{thm:conclusion-multiscale-stokes-energy-l2-reconstruction} 的定义，\(\mathcal E_m(\alpha)\) 正是 \(d\alpha\) 在尺度 \(m\) 的 dyadic 胞腔平均所形成的分片常系数 \((k+1)\)-形式的 \(L^2\) 范数平方。将 \(d\alpha\) 写成系数向量 \((f_I)_{|I|=k+1}\) 后，这恰好是按 \(\mathscr F_m\) 作条件期望所得的分片常值形式，因此
\[
\mathcal E_m(\alpha)
=
\sum_{|I|=k+1}\|P_m f_I\|_2^2
=
\|P_m(d\alpha)\|_2^2.
\]
再令
\[
D_m:=P_{m+1}(d\alpha)-P_m(d\alpha).
\]
由于 \((P_m(d\alpha))_{m\ge 0}\) 为 \(L^2\) 鞅，\(D_m\) 与 \(P_m(d\alpha)\) 在 \(L^2\) 中正交，故
\[
\mathcal E_{m+1}(\alpha)-\mathcal E_m(\alpha)
=
\|P_{m+1}(d\alpha)\|_2^2-\|P_m(d\alpha)\|_2^2
=
\|D_m\|_2^2.
\]
为得到局部公式，只需在每个 \(Q\in\mathcal Q_m\) 上计算。由于 \(P_m(d\alpha)\) 在 \(Q\) 上为常值，而 \(P_{m+1}(d\alpha)\) 在 \(Q\) 的每个子立方体 \(Q'\prec Q\) 上亦为常值，且父平均等于子平均的算术平均，
\[
P_m(d\alpha)|_Q
=
\frac1{2^d}\sum_{Q'\prec Q}P_{m+1}(d\alpha)|_{Q'}.
\]
因此对固定 \(Q\) 应用有限维恒等式
\[
\frac1{2^d}\sum_{i=1}^{2^d}|u_i|^2
-\left|\frac1{2^d}\sum_{i=1}^{2^d}u_i\right|^2
=
\frac1{2^d}\sum_{i=1}^{2^d}|u_i-\bar u|^2,
\qquad
\bar u:=\frac1{2^d}\sum_i u_i,
\]
并令 \(u_i=P_{m+1}(d\alpha)|_{Q_i'}\)，即可得到
\[
\int_Q |D_m|^2
=
\frac{|Q|}{2^d}\sum_{Q'\prec Q}
\bigl|
P_{m+1}(d\alpha)|_{Q'}-P_m(d\alpha)|_Q
\bigr|^2.
\]
对所有 \(Q\in\mathcal Q_m\) 求和，并注意 \(|Q|=h_m^d\)，便得到所示显式增量公式。右端非负，故单调性随之成立。证毕。
\end{proof}

\begin{corollary}[尺度增量谱的后验误差与 dyadic Besov 读出]\label{cor:conclusion-stokes-energy-dyadic-besov-readout}
在定理 \ref{thm:conclusion-stokes-energy-general-degree-cascade} 的设定下，对任意 \(m\ge 0\) 有
\[
\|d\alpha-P_m(d\alpha)\|_{L^2(Q)}^2
=
\|d\alpha\|_{L^2(Q)}^2-\mathcal E_m(\alpha)
=
\sum_{j=m}^{\infty}\bigl(\mathcal E_{j+1}(\alpha)-\mathcal E_j(\alpha)\bigr).
\]
因此右端给出了 \(\|d\alpha-P_m(d\alpha)\|_{L^2}\) 的纯尺度后验误差证书。

进一步，对任意 \(s\in(0,1)\) 定义
\[
\mathcal N_s(\alpha)^2
:=
\sum_{m\ge 0}2^{2sm}\bigl(\mathcal E_{m+1}(\alpha)-\mathcal E_m(\alpha)\bigr).
\]
则存在仅依赖于 \((d,k,s)\) 的常数 \(0<c_{d,k,s}\le C_{d,k,s}<\infty\)，使
\[
c_{d,k,s}\,\mathcal N_s(\alpha)^2
\le
\sum_{|I|=k+1}\sum_{m\ge 0}2^{2sm}\|P_{m+1}f_I-P_mf_I\|_{L^2(Q)}^2
\le
C_{d,k,s}\,\mathcal N_s(\alpha)^2.
\]
特别地，\(\mathcal N_s(\alpha)\) 与 \(d\alpha\) 的 dyadic Besov 半范数等价，故其完全由多尺度 Stokes 增量数据决定。
\end{corollary}
\begin{proof}
由 \(L^2(Q;\Lambda^{k+1})\) 中的鞅收敛定理，
\[
d\alpha
=
P_m(d\alpha)+\sum_{j=m}^{\infty}\bigl(P_{j+1}(d\alpha)-P_j(d\alpha)\bigr),
\]
且各级差分两两正交。于是
\[
\|d\alpha-P_m(d\alpha)\|_2^2
=
\sum_{j=m}^{\infty}\|P_{j+1}(d\alpha)-P_j(d\alpha)\|_2^2.
\]
再由定理 \ref{thm:conclusion-stokes-energy-general-degree-cascade}，
\[
\|P_{j+1}(d\alpha)-P_j(d\alpha)\|_2^2
=
\mathcal E_{j+1}(\alpha)-\mathcal E_j(\alpha),
\]
从而得到第一组恒等式。

对第二条，只需把 \(d\alpha\) 的各个系数分量 \(f_I\) 视作标量 \(L^2\) 函数，并应用 dyadic Haar 基下的标准 Littlewood--Paley/Besov 刻画。由于
\[
\|P_{m+1}(d\alpha)-P_m(d\alpha)\|_2^2
=
\sum_{|I|=k+1}\|P_{m+1}f_I-P_mf_I\|_2^2,
\]
加权求和便给出所述等价性。证毕。
\end{proof}

\endinput


以上讨论在本质上依赖于可交换情形下的线性斯托克斯公式与能量可加性。
当边界量由非阿贝尔平行移动给出时，“曲率通量主导、三阶余项控制”的局部结构仍可在小尺度上保持，并可通过 dyadic 递归实现全域误差控制。

\begin{theorem}[非阿贝尔情形下的小矩形广义斯托克斯定量化]\label{thm:conclusion-nonabelian-stokes-small-rectangle}
\leanverified{paper\_conclusion\_nonabelian\_stokes\_small\_rectangle}
设 \(G\subset \mathrm{GL}_N(\mathbb C)\) 为矩阵李群，\(\mathfrak g\) 为其李代数，\(\|\cdot\|\) 为算子范数。
令
\[
A:=A_1\,\mathrm d x_1+A_2\,\mathrm d x_2
\]
为定义在包含 \(R_h:=[0,h]\times[0,h]\) 的开邻域上的 \(\mathfrak g\)-值 \(C^2\) \(1\)-形式，并记曲率
\[
F_A
:=\Bigl(\partial_1A_2-\partial_2A_1+[A_1,A_2]\Bigr)\,\mathrm d x_1\wedge \mathrm d x_2
=:F_{12}\,\mathrm d x_1\wedge \mathrm d x_2.
\]
令 \(\operatorname{Hol}_A(\partial R_h)\in G\) 表示沿 \(\partial R_h\)（逆时针取向）的平行移动回路整体单值。
则存在绝对常数 \(C>0\) 使得：若
\[
h\,\|A\|_{L^\infty(R_h)}\le \frac{1}{10},
\]
则 \(\operatorname{Hol}_A(\partial R_h)\) 落入 \(\exp\) 的局部可逆邻域内，从而存在唯一 \(\Omega_h\in\mathfrak g\) 使
\[
\operatorname{Hol}_A(\partial R_h)=\exp(\Omega_h),
\]
并且满足定量误差估计
\[
\Bigl\|\Omega_h-\iint_{R_h}F_{12}(x)\,\mathrm d x_1\,\mathrm d x_2\Bigr\|
\le
C\,h^3\exp\!\bigl(C h\|A\|_{L^\infty(R_h)}\bigr)
\Bigl(
\|\nabla F_{12}\|_{L^\infty(R_h)}+\|A\|_{L^\infty(R_h)}\|F_{12}\|_{L^\infty(R_h)}
\Bigr),
\]
其中 \(\nabla F_{12}\) 表示 \(F_{12}\) 的一阶导数张量。
若 \(\mathfrak g\) 交换，则上式右端为 \(0\)，从而严格有
\[
\Omega_h=\iint_{R_h}F_{12}\,\mathrm d x_1\,\mathrm d x_2.
\]
\end{theorem}
\begin{proof}[证明要点]
将 \(\operatorname{Hol}_A(\partial R_h)\) 写为四条边上传输算子的有序乘积，并对每条边的路径有序指数作二阶截断（等价地，可采用 Magnus 展开的一阶项与余项估计）。
在 \(\|A\|_{L^\infty}\) 受控下，二阶主项由线性斯托克斯给出 \(\iint(\partial_1A_2-\partial_2A_1)\)，而非交换性导致的二阶交叉项经 BCH 归并产生 \(\iint[A_1,A_2]\)，从而得到 \(\iint F_{12}\)。
三阶余项由 \(C^2\) 正则性给出，并可用 \(\|\nabla F_{12}\|_{L^\infty}\) 与共轭扰动项 \(\|A\|_{L^\infty}\|F_{12}\|_{L^\infty}\) 控制，同时携带 \(\exp(C h\|A\|_{L^\infty})\) 放大因子。
在 \(h\|A\|_{L^\infty}\le 1/10\) 下，\(\log\) 在单位元邻域内为 Lipschitz，从而把对 \(\operatorname{Hol}_A(\partial R_h)\) 的误差估计传递到 \(\Omega_h\)。证毕。
\end{proof}

\begin{corollary}[dyadic 网格上的递归斯托克斯证书与多尺度误差控制]\label{cor:conclusion-nonabelian-dyadic-stokes-certificate}
在定理 \ref{thm:conclusion-nonabelian-stokes-small-rectangle} 的设定下，令 \(R_1=[0,1]^2\) 并作 dyadic 细分为边长 \(h:=2^{-m}\) 的小方形集合 \(\mathcal Q_m\)。
对每个 \(Q\in\mathcal Q_m\)，设其边界 holonomy 的对数存在且记为 \(\Omega(Q)\in\mathfrak g\)（当 \(m\) 充分大时由定理 \ref{thm:conclusion-nonabelian-stokes-small-rectangle} 保证）。
则存在绝对常数 \(C>0\)（与 \(m\) 无关）使得当 \(m\) 充分大时
\[
\Bigl\|\sum_{Q\in\mathcal Q_m}\Omega(Q)-\iint_{R_1}F_{12}\,\mathrm d x_1\,\mathrm d x_2\Bigr\|
\le
C\,2^{-m}\exp\!\bigl(C\|A\|_{L^\infty(R_1)}\bigr)
\Bigl(
\|\nabla F_{12}\|_{L^\infty(R_1)}+\|A\|_{L^\infty(R_1)}\|F_{12}\|_{L^\infty(R_1)}
\Bigr).
\]
\end{corollary}
\begin{proof}
对每个 \(Q\in\mathcal Q_m\) 应用定理 \ref{thm:conclusion-nonabelian-stokes-small-rectangle} 得
\(
\Omega(Q)=\iint_QF_{12}+E(Q)
\)
且 \(\|E(Q)\|\le C h^3\exp(C h\|A\|_{L^\infty(R_1)})(\cdots)\)。
对全部小方形求和，主项严格相加为 \(\iint_{R_1}F_{12}\)，而误差满足
\[
\sum_{Q\in\mathcal Q_m}\|E(Q)\|
\le
|\mathcal Q_m|\cdot C h^3\exp\!\bigl(C\|A\|_{L^\infty(R_1)}\bigr)(\cdots)
=C h\,\exp\!\bigl(C\|A\|_{L^\infty(R_1)}\bigr)(\cdots),
\]
即得到所示界。证毕。
\end{proof}

\paragraph{非阿贝尔 dyadic holonomy 的曲率能量与缺陷级联}
二维小回路的对数 holonomy 估计可在 dyadic 细化下闭合为一条全局 \(L^2\) 曲率重建定理；与此同时，单块量化误差能够完全重写为一列可望远镜的通量缺陷，从而给出纯边界数据驱动的后验证书。

\begin{theorem}[非阿贝尔多尺度 Stokes 能量的曲率重建与一阶量化误差]\label{thm:conclusion-nonabelian-dyadic-stokes-energy-reconstruction}
\leanverified{paper\_conclusion\_nonabelian\_dyadic\_stokes\_energy\_reconstruction}
沿用定理 \ref{thm:conclusion-nonabelian-stokes-small-rectangle} 的设定，并令 \(R_1=[0,1]^2\)。
记 \(\mathcal Q_m\) 为 \(R_1\) 的 dyadic 方形剖分，边长 \(h_m:=2^{-m}\)。
当 \(m\) 充分大时，对每个 \(Q\in\mathcal Q_m\) 都有 \(h_m\|A\|_{L^\infty(Q)}\le 1/10\)，故存在唯一 \(\Omega(Q)\in\mathfrak g\) 满足
\[
\operatorname{Hol}_A(\partial Q)=\exp(\Omega(Q)).
\]
定义非阿贝尔多尺度 Stokes 能量
\[
\mathcal E_m^{\mathrm{NA}}(A)
:=
h_m^2\sum_{Q\in\mathcal Q_m}
\left\|
\frac{\Omega(Q)}{h_m^2}
\right\|^2.
\]
则有收敛
\[
\mathcal E_m^{\mathrm{NA}}(A)\longrightarrow \int_{R_1}|F_{12}(x)|^2\,\mathrm d x
\qquad (m\to\infty).
\]
若进一步 \(F_{12}\) 为 Lipschitz，则存在仅依赖于群嵌入与维数的常数 \(C>0\) 使对充分大的 \(m\) 有
\[
\left|
\mathcal E_m^{\mathrm{NA}}(A)-\int_{R_1}|F_{12}|^2
\right|
\le
C h_m e^{C\|A\|_{L^\infty(R_1)}}
\Bigl(
\|F_{12}\|_{L^\infty(R_1)}\|\nabla F_{12}\|_{L^\infty(R_1)}
+\|A\|_{L^\infty(R_1)}\|F_{12}\|_{L^\infty(R_1)}^2
+\|\nabla F_{12}\|_{L^\infty(R_1)}^2
\Bigr).
\]
\end{theorem}
\begin{proof}
对每个 \(Q\in\mathcal Q_m\) 记 dyadic 曲率平均
\[
\overline F_Q
:=
\frac1{h_m^2}\int_Q F_{12}(x)\,\mathrm d x.
\]
由定理 \ref{thm:conclusion-nonabelian-stokes-small-rectangle}，存在常数 \(C>0\) 使
\[
\left\|
\frac{\Omega(Q)}{h_m^2}-\overline F_Q
\right\|
\le
C h_m e^{C\|A\|_{L^\infty(R_1)}}
\Bigl(
\|\nabla F_{12}\|_{L^\infty(R_1)}
+\|A\|_{L^\infty(R_1)}\|F_{12}\|_{L^\infty(R_1)}
\Bigr)
\]
对所有 \(Q\in\mathcal Q_m\) 同时成立。
令
\[
\mathscr H_m(x)
:=
\sum_{Q\in\mathcal Q_m}\frac{\Omega(Q)}{h_m^2}\,\mathbf 1_Q(x),
\qquad
P_mF_{12}(x)
:=
\sum_{Q\in\mathcal Q_m}\overline F_Q\,\mathbf 1_Q(x).
\]
则
\[
\mathcal E_m^{\mathrm{NA}}(A)=\|\mathscr H_m\|_{L^2(R_1)}^2,
\qquad
\|P_mF_{12}\|_{L^2(R_1)}^2
=
h_m^2\sum_{Q\in\mathcal Q_m}|\overline F_Q|^2.
\]
并且上面的点态估计给出
\[
\|\mathscr H_m-P_mF_{12}\|_{L^\infty(R_1)}
\le
C h_m e^{C\|A\|_{L^\infty(R_1)}}
\Bigl(
\|\nabla F_{12}\|_{L^\infty(R_1)}
+\|A\|_{L^\infty(R_1)}\|F_{12}\|_{L^\infty(R_1)}
\Bigr).
\]
由于 \(|R_1|=1\)，同一估计也成立于 \(L^2\) 范数。于是
\[
\left|
\mathcal E_m^{\mathrm{NA}}(A)-\|P_mF_{12}\|_2^2
\right|
\le
\|\mathscr H_m-P_mF_{12}\|_2\,
\bigl(\|\mathscr H_m\|_2+\|P_mF_{12}\|_2\bigr).
\]
再由 \(\|P_mF_{12}\|_2\le \|F_{12}\|_\infty\) 与
\[
\|\mathscr H_m\|_2
\le
\|P_mF_{12}\|_2+\|\mathscr H_m-P_mF_{12}\|_2
\]
可得
\[
\left|
\mathcal E_m^{\mathrm{NA}}(A)-\|P_mF_{12}\|_2^2
\right|
\le
C h_m e^{C\|A\|_\infty}
\Bigl(
\|F_{12}\|_\infty\|\nabla F_{12}\|_\infty
+\|A\|_\infty\|F_{12}\|_\infty^2
+\|\nabla F_{12}\|_\infty^2
\Bigr),
\]
其中把二次项中可能出现的 \(\|A\|_\infty^2\|F_{12}\|_\infty^2\) 吸收到指数因子中。
另一方面，\(P_mF_{12}\) 正是 \(F_{12}\) 相对于 dyadic \(\sigma\)-代数的条件期望，因此 \(P_mF_{12}\to F_{12}\) 于 \(L^2(R_1)\)。故
\[
\|P_mF_{12}\|_2^2\longrightarrow \|F_{12}\|_2^2
=
\int_{R_1}|F_{12}(x)|^2\,\mathrm d x,
\]
与上式合并即得收敛及一阶误差界。证毕。
\end{proof}

\begin{theorem}[dyadic 通量缺陷的完全望远镜分解]\label{thm:conclusion-nonabelian-dyadic-flux-defect-telescoping}
\leanverified{paper\_conclusion\_nonabelian\_dyadic\_flux\_defect\_telescoping}
在定理 \ref{thm:conclusion-nonabelian-dyadic-stokes-energy-reconstruction} 的设定下，固定某个 dyadic 方形 \(Q_0\subset R_1\)，设其边长为 \(h_0=2^{-m_0}\)。
对 \(r\ge 0\) 记 \(\mathcal Q_{m_0+r}(Q_0)\) 为 \(Q_0\) 的 \(r\) 次 dyadic 细分，并定义
\[
S_r(Q_0):=\sum_{Q\in\mathcal Q_{m_0+r}(Q_0)}\Omega(Q)\in\mathfrak g.
\]
则 \(S_r(Q_0)\) 在 \(\mathfrak g\) 中收敛，且
\[
\lim_{r\to\infty}S_r(Q_0)=\int_{Q_0}F_{12}(x)\,\mathrm d x.
\]
进一步，对每个层级 \(j\ge m_0\) 与 \(Q\in\mathcal Q_j\) 定义 dyadic 通量缺陷
\[
\Delta_j(Q):=\Omega(Q)-\sum_{Q'\prec Q}\Omega(Q').
\]
则级数
\[
\sum_{j=m_0}^{\infty}\ \sum_{\substack{Q\in\mathcal Q_j\\ Q\subseteq Q_0}}\Delta_j(Q)
\]
绝对收敛，且成立完全望远镜恒等式
\[
\Omega(Q_0)-\int_{Q_0}F_{12}(x)\,\mathrm d x
=
\sum_{j=m_0}^{\infty}\ \sum_{\substack{Q\in\mathcal Q_j\\ Q\subseteq Q_0}}\Delta_j(Q).
\]
因此得到纯边界数据型误差证书
\[
\left\|
\Omega(Q_0)-\int_{Q_0}F_{12}(x)\,\mathrm d x
\right\|
\le
\sum_{j=m_0}^{\infty}\ \sum_{\substack{Q\in\mathcal Q_j\\ Q\subseteq Q_0}}
\|\Delta_j(Q)\|.
\]
\end{theorem}
\begin{proof}
对任意 dyadic 方形 \(Q\subset R_1\) 记
\[
E(Q):=\Omega(Q)-\int_Q F_{12}(x)\,\mathrm d x.
\]
由定理 \ref{thm:conclusion-nonabelian-stokes-small-rectangle}，存在常数
\[
M_A
:=
e^{C\|A\|_{L^\infty(R_1)}}
\Bigl(
\|\nabla F_{12}\|_{L^\infty(R_1)}
+\|A\|_{L^\infty(R_1)}\|F_{12}\|_{L^\infty(R_1)}
\Bigr)
\]
使对任意边长为 \(h\) 的 dyadic 方形 \(Q\) 有
\[
\|E(Q)\|\le C h^3 M_A.
\]
于是对边长 \(h_r:=2^{-(m_0+r)}\) 的细分层，
\[
S_r(Q_0)
=
\int_{Q_0}F_{12}(x)\,\mathrm d x
+\sum_{Q\in\mathcal Q_{m_0+r}(Q_0)}E(Q).
\]
由于 \(|\mathcal Q_{m_0+r}(Q_0)|=4^r\)，故
\[
\sum_{Q\in\mathcal Q_{m_0+r}(Q_0)}\|E(Q)\|
\le
4^r\cdot C h_r^3 M_A
=
C h_0^2 2^{-r} M_A\longrightarrow 0,
\]
从而 \(S_r(Q_0)\to \int_{Q_0}F_{12}\)。

再来看通量缺陷。由积分在子方块上的可加性，
\[
\Delta_j(Q)
=
\Bigl(\Omega(Q)-\int_Q F_{12}\Bigr)
-\sum_{Q'\prec Q}\Bigl(\Omega(Q')-\int_{Q'}F_{12}\Bigr)
=
E(Q)-\sum_{Q'\prec Q}E(Q').
\]
因此对边长 \(2^{-j}\) 的 \(Q\) 有
\[
\|\Delta_j(Q)\|
\le
\|E(Q)\|+\sum_{Q'\prec Q}\|E(Q')\|
\le
C 2^{-3j}M_A.
\]
而 \(Q_0\) 内第 \(j\) 层方形数为 \(4^{j-m_0}\)，故
\[
\sum_{\substack{Q\in\mathcal Q_j\\ Q\subseteq Q_0}}\|\Delta_j(Q)\|
\le
C 4^{j-m_0}2^{-3j}M_A
=
C 2^{-2m_0-j}M_A.
\]
右端对 \(j\) 可求和，故双重级数绝对收敛。

最后，由定义
\[
\Omega(Q)=\sum_{Q'\prec Q}\Omega(Q')+\Delta_j(Q).
\]
对 \(Q_0\) 内第 \(j\) 层所有方形求和，得到
\[
S_{j-m_0}(Q_0)
=
S_{j+1-m_0}(Q_0)
+\sum_{\substack{Q\in\mathcal Q_j\\ Q\subseteq Q_0}}\Delta_j(Q).
\]
从 \(j=m_0\) 望远镜求和到 \(m_0+r-1\) 即得
\[
\Omega(Q_0)-S_r(Q_0)
=
\sum_{j=m_0}^{m_0+r-1}\ \sum_{\substack{Q\in\mathcal Q_j\\ Q\subseteq Q_0}}\Delta_j(Q).
\]
令 \(r\to\infty\)，利用前面已证 \(S_r(Q_0)\to \int_{Q_0}F_{12}\) 及绝对收敛性，即得完全望远镜恒等式。最后对两边取范数并用三角不等式便得到误差证书。证毕。
\end{proof}

\begin{conjecture}[holonomy--dyadic 缺陷的非阿贝尔 Littlewood--Paley 框架]\label{conj:conclusion-nonabelian-dyadic-littlewood-paley}
沿用定理 \ref{thm:conclusion-nonabelian-dyadic-flux-defect-telescoping} 的记号，并假设存在全局规范使所有 dyadic 层的 \(\log\circ \operatorname{Hol}_A\) 都在同一主支上定义。
令 \(h_m:=2^{-m}\)。对 \(s\in(0,1)\) 定义加权平方函数
\[
\mathfrak L_s(A)^2
:=
\sum_{m\ge 0}
2^{2sm} h_m^2
\sum_{Q\in\mathcal Q_m}
\left\|
\frac{\Delta_m(Q)}{h_m^2}
\right\|^2.
\]
则预期存在仅依赖于群嵌入与 \(s\) 的常数 \(0<c_s\le C_s<\infty\)，使
\[
c_s\,|F_{12}|_{\dot B_{2,2}^{s,\mathrm{dy}}(R_1)}^2
\le
\mathfrak L_s(A)^2
\le
C_s\,|F_{12}|_{\dot B_{2,2}^{s,\mathrm{dy}}(R_1)}^2.
\]
当 \(s=0\) 时，上式应退化为与杨--米尔斯作用量
\[
\int_{R_1}|F_{12}(x)|^2\,\mathrm d x
\]
等价的 dyadic square-function 读出；与定理 \ref{thm:conclusion-nonabelian-dyadic-flux-defect-telescoping} 结合后，这将把非阿贝尔曲率能量、正则度与后验误差统一到同一条 holonomy 驱动的多尺度证书链中。
\end{conjecture}

\endinput


\begin{theorem}[三维非交换 Stokes--Bianchi 立方体证书与四阶误差界]\label{thm:conclusion-nonabelian-stokes-bianchi-cube-certificate}
设 \(G\subset \mathrm{GL}_N(\mathbb C)\) 为矩阵李群，\(\mathfrak g\) 为其李代数，\(\|\cdot\|\) 为算子范数。
令
\[
A:=A_1\,\mathrm d x_1+A_2\,\mathrm d x_2+A_3\,\mathrm d x_3
\]
为定义在包含闭立方体 \([x,x+h]^3\) 的开邻域上的 \(\mathfrak g\)-值 \(C^3\) \(1\)-形式（\(h>0\) 足够小），并记曲率 \(F:=\mathrm d A+A\wedge A\) 与协变导数 \(\nabla_A F\)。
对 \(i=1,2,3\) 定义沿坐标边的平行移动
\[
U_i(x):=\mathcal P\exp\!\left(\int_0^h A_i(x+t e_i)\,\mathrm d t\right)\in G,
\]
并对 \(1\le i<j\le 3\) 定义基点在 \(x\) 的面元 holonomy
\[
U_{ij}(x):=U_i(x)\,U_j(x+h e_i)\,U_i(x+h e_j)^{-1}\,U_j(x)^{-1}\in G.
\]
再定义立方体 Bianchi 缺陷
\begin{equation}\label{eq:conclusion-bianchi-defect-cube}
\begin{aligned}
B_h(x):={}&U_{12}(x)\,U_{23}(x)\,U_{31}(x)\\
&\cdot\bigl(U_3(x)\,U_{12}(x+h e_3)\,U_3(x)^{-1}\bigr)^{-1}
\cdot\bigl(U_1(x)\,U_{23}(x+h e_1)\,U_1(x)^{-1}\bigr)^{-1}\\
&\cdot\bigl(U_2(x)\,U_{31}(x+h e_2)\,U_2(x)^{-1}\bigr)^{-1}.
\end{aligned}
\end{equation}
则存在仅依赖于 \(N\) 的常数 \(C>0\) 使得当 \(h\,\|A\|_{L^\infty([x,x+h]^3)}\le 1/10\) 时，\(B_h(x)\) 落入 \(\exp\) 的局部可逆邻域内，从而主支对数 \(\log B_h(x)\in\mathfrak g\) 定义且满足估计
\begin{equation}\label{eq:conclusion-bianchi-defect-order4}
\boxed{\;
\|\log B_h(x)\|
\le
C\,h^4\Bigl(\|F\|_{C^1([x,x+h]^3)}^{2}+\|\nabla_A F\|_{C^0([x,x+h]^3)}\Bigr)\; }.
\end{equation}
更强地，存在三阶主导展开
\begin{equation}\label{eq:conclusion-bianchi-defect-leading}
\boxed{\;
\log B_h(x)
=h^3\bigl(\dd_A F\bigr)_{123}(x)
O\!\left(h^4\Bigl(\|F\|_{C^1([x,x+h]^3)}^{2}+\|\nabla_A F\|_{C^0([x,x+h]^3)}\Bigr)\right)\; }.
\end{equation}
由于 Bianchi 恒等式 \(\dd_A F\equiv 0\) 对任意光滑连接成立，\eqref{eq:conclusion-bianchi-defect-leading} 立即推出 \eqref{eq:conclusion-bianchi-defect-order4} 的四阶衰减：\(B_h(x)\) 给出“协变闭性”的可审计立方体证书。
\end{theorem}
\begin{proof}[证明要点]
对每个面元 holonomy \(U_{ij}(x)\) 用路径有序指数的 Dyson/Magnus 展开得到二阶主项
\(
U_{ij}(x)=I+h^2F_{ij}(x)+O(h^3)
\)。
将对顶面元搬运回同一基点所产生的共轭项展开到三阶，可把 \(\partial_kF_{ij}+[A_k,F_{ij}]\) 归并为 \((\nabla_{A,k}F_{ij})(x)\)。
把这些展开代入 \eqref{eq:conclusion-bianchi-defect-cube}，二阶曲率项在乘积中完全抵消，三阶项收束为
\(
\nabla_{A,3}F_{12}+\nabla_{A,1}F_{23}+\nabla_{A,2}F_{31}
=(\dd_A F)_{123}
\)。
其余项由 BCH 公式与三阶展开的余项估计控制；非交换修正从交换子起始，至少为 \(O(h^4\|F\|^2)\)，并与 \(\|\nabla_A F\|_{C^0}\) 一并被 \eqref{eq:conclusion-bianchi-defect-leading} 右端吸收。
在 \(h\|A\|_{L^\infty}\le 1/10\) 下，\(\log\) 在单位元邻域内 Lipschitz，从而得到 \eqref{eq:conclusion-bianchi-defect-order4}。证毕。
\end{proof}

\begin{corollary}[dyadic 三维网格上的 Bianchi 审计阈值]\label{cor:conclusion-nonabelian-bianchi-dyadic-audit-threshold}
在定理 \ref{thm:conclusion-nonabelian-stokes-bianchi-cube-certificate} 的设定下，令 \(h_m:=2^{-m}\) 并定义
\[
\mathcal B_m
:=\sup_{x\in h_m\ZZ^{3}\cap[0,1-h_m]^3}\|\log B_{h_m}(x)\|.
\]
则存在常数 \(C>0\) 使当 \(m\) 充分大时
\[
\boxed{\;
\mathcal B_m
\le
C\,2^{-4m}\Bigl(\|F\|_{C^1([0,1]^3)}^{2}+\|\nabla_A F\|_{C^0([0,1]^3)}\Bigr)\; }.
\]
因此若某一数值实现对大 \(m\) 的观测衰减指数严格小于 \(4\)，则该实现与协变 Stokes--Bianchi 结构不相容，并构成离散化破坏协变性的直接否证证书。
\end{corollary}
\begin{proof}
对每个 dyadic 小立方体 \([x,x+h_m]^3\) 应用 \eqref{eq:conclusion-bianchi-defect-order4} 并取上确界即可。证毕。
\end{proof}

\begin{theorem}[相对斯托克斯与圆维缺陷：由边界积分实现的同调维数刻画]\label{thm:conclusion-relative-stokes-torus-defect}
\leanverified{paper\_conclusion\_relative\_stokes\_torus\_defect}
设 \(f:G\to H\) 为有限生成阿贝尔群同态，记 \(K:=\ker f\)、\(I:=\operatorname{im}f\)。
令 \(\widehat{(\cdot)}\) 表 Pontryagin 对偶，并令 \(T_X\) 表示紧交换群 \(\widehat X\) 的恒等分量（故 \(T_X\) 为某一实环面）。
定义
\[
\delta(f):=\operatorname{rank}K.
\]
则：
\begin{enumerate}
  \item 由短正合列 \(0\to K\to G\to I\to 0\) 诱导出环面短正合列
  \[
  0\longrightarrow T_I\longrightarrow T_G\longrightarrow T_K\longrightarrow 0,
  \qquad
  T_K\simeq (S^1)^{\delta(f)}.
  \]
  \item 存在自然同构
  \[
  H^1_{\mathrm{dR}}(T_G,T_I;\mathbb R)\ \cong\ H^1_{\mathrm{dR}}(T_K;\mathbb R),
  \]
  从而
  \[
  \dim_{\mathbb R}H^1_{\mathrm{dR}}(T_G,T_I;\mathbb R)=\delta(f).
  \]
  \item 在整周期晶格意义下有自然同构
  \[
  H^1_{\mathrm{dR}}(T_G,T_I;\mathbb Z)\ \cong\ \operatorname{Hom}(K,\mathbb Z).
  \]
  特别地，对任意相对 \(1\)-循环 \(\gamma\in H_1(T_G,T_I;\mathbb Z)\) 与任意相对闭 \(1\)-形式
  \([\omega]\in H^1_{\mathrm{dR}}(T_G,T_I;\mathbb Z)\)，积分配对
  \(
  \langle[\omega],\gamma\rangle:=\int_\gamma\omega
  \)
  仅依赖于相对同调类。
\end{enumerate}
\end{theorem}
\begin{proof}[证明要点]
有限生成阿贝尔群分解 \(X\simeq \mathbb Z^{r_X}\oplus \mathrm{Tor}(X)\)；其对偶的恒等分量为 \(T_X\simeq (S^1)^{r_X}\)，从而 \(r_K=\delta(f)\) 给出 (1)。
又因 \(T_I\subset T_G\) 为闭子环面且商 \(T_G/T_I\simeq T_K\)，相对 de Rham 上同调与商空间的约化上同调在该类情形下自然同构，得到 (2)。
整周期部分由
\(
H^1_{\mathrm{dR}}((S^1)^n;\mathbb Z)\cong \operatorname{Hom}(\mathbb Z^n,\mathbb Z)
\)
与自由部分识别得到 (3)，而积分配对的同调不变性即相对斯托克斯公式。证毕。
\end{proof}

\begin{proposition}[离散广义斯托克斯的可积化校正]\label{prop:conclusion-discrete-stokes-integrable-correction}
\leanverified{paper\_conclusion\_discrete\_stokes\_integrable\_correction}
令 \(X\) 为有限、连通的立方体复形，记 \(C^k(X;\RR)\) 为 \(k\)-余链群，\(\delta:C^k\to C^{k+1}\) 为余边界算子。
对任意 \(k\ge 1\) 与任意 \(\beta\in C^k(X;\RR)\)，存在 \(\alpha\in C^{k-1}(X;\RR)\) 使得残差
\[
r:=\beta-\delta\alpha
\]
满足统一控制
\[
\|r\|_{\ell^\infty(\mathcal C_k)}\ \le\ \mathrm{diam}(X)\,\|\delta\beta\|_{\ell^\infty(\mathcal C_{k+1})},
\]
其中 \(\mathcal C_j\) 表示 \(X\) 的定向 \(j\)-胞腔集合，\(\mathrm{diam}(X)\) 为 \(1\)-骨架的图距离直径。
并且当 \(\delta\beta=0\) 时可取 \(r=0\)，从而在可收缩立方体复形上闭 \(k\)-余链必为恰当余链。
\end{proposition}
\begin{proof}[证明要点]
取 \(1\)-骨架上一棵生成树 \(T\) 并固定根顶点 \(v_0\)。
对每个定向 \(k\)-胞腔 \(\sigma\)，选取其某个顶点 \(v(\sigma)\)，并令 \(\gamma_{v(\sigma)}\) 为 \(T\) 中从 \(v_0\) 到 \(v(\sigma)\) 的唯一路径；其长度不超过 \(\mathrm{diam}(X)\)。
在立方体复形中可构造链同伦锥算子 \(\mathrm{Cone}:C_k(X)\to C_{k+1}(X)\)，使
\[
\partial\,\mathrm{Cone}+\mathrm{Cone}\,\partial=\mathrm{id}_{C_k(X)}\qquad(k\ge 1),
\]
且对每个 \(\sigma\) 的锥链由至多 \(\mathrm{diam}(X)\) 个 \((k+1)\)-胞腔组成。
对偶化令 \(H:=\mathrm{Cone}^\ast:C^k(X)\to C^{k-1}(X)\)，则
\[
\delta H+H\delta=\mathrm{id}_{C^k(X)}\qquad(k\ge 1),
\qquad
\|H\|_{\ell^\infty\to\ell^\infty}\le \mathrm{diam}(X).
\]
取 \(\alpha:=H\beta\)，则
\(
r=\beta-\delta\alpha=\beta-\delta H\beta=H\delta\beta
\)，从而
\(
\|r\|_\infty\le \|H\|\,\|\delta\beta\|_\infty\le \mathrm{diam}(X)\|\delta\beta\|_\infty
\)。
当 \(\delta\beta=0\) 时得 \(r=0\)。证毕。
\end{proof}

\begin{theorem}[均匀网格上的量化 Stokes--Poincar\'e 分解与尺度常数]\label{thm:conclusion-uniform-grid-quantized-stokes-poincare-scale}
令 \(I^n=[0,1]^n\)，并取网格尺度 \(h=2^{-m}\) 的均匀立方剖分 \(\mathcal Q_h\)。
对 \(k\)-余链 \(\beta\in C^k(\mathcal Q_h;\RR)\) 定义密度范数
\[
\|\beta\|_{\infty,h}:=\max_{\sigma^k\in\mathcal Q_h}\frac{|\beta(\sigma^k)|}{\mathrm{vol}(\sigma^k)}
\qquad(\mathrm{vol}(\sigma^k)=h^k).
\]
则对任意 \(k\ge 1\) 存在线性算子
\[
H_{k,h}:C^k(\mathcal Q_h;\RR)\longrightarrow C^{k-1}(\mathcal Q_h;\RR)
\]
满足同伦恒等式
\[
\delta H_{k,h}+H_{k+1,h}\delta=\mathrm{id}_{C^k(\mathcal Q_h)}
\]
以及尺度最优的算子范数控制
\[
\|H_{k,h}\|_{(\infty,h)\to(\infty,h)}\ \le\ \frac{h}{2k}.
\]
从而对任意 \(\beta\in C^k(\mathcal Q_h;\RR)\)，取 \(\alpha:=H_{k,h}\beta\) 有
\[
\|\alpha\|_{\infty,h}\le \frac{h}{2k}\|\beta\|_{\infty,h},
\qquad
\|\beta-\delta\alpha\|_{\infty,h}\le \frac{h}{2k}\|\delta\beta\|_{\infty,h}.
\]
\end{theorem}
\begin{proof}[证明要点]
对单位立方复形的余链复形，定理 \ref{thm:spg-cubical-quantized-stokes-poincare-linf-optimal} 给出同伦算子与 \(\ell^\infty\) 最优常数 \(\frac{1}{2k}\)。
将任一边长 \(h\) 的小立方体通过仿射缩放同构到单位立方体，\(\delta\) 与同伦恒等式在该同构下保持不变，而密度范数在缩放下按长度线性缩放，从而常数相应乘以 \(h\)。
在全体小立方体上以一致方式拼接该局部算子即可得到 \(H_{k,h}\)；算子范数与同伦恒等式由局部构造与立方剖分的兼容性给出。证毕。
\end{proof}

\begin{theorem}[\(\frac14\) 常数的跨范畴刚性]\label{thm:conclusion-onequarter-cross-category-rigidity}
常数
\[
\frac14
\]
并非某个单独模型中的数值巧合，而是同一类一阶 Stokes 正规化在多层实现中的 sharp 常数。更准确地说，在下列四种彼此兼容的实现范畴中，
\[
\text{光滑 de Rham 同伦}
\longrightarrow
\text{立方余链}
\longrightarrow
\text{超立方体边流}
\longrightarrow
\text{连续时间跳跃链},
\]
凡把二维方形缺陷正规化为一维势差残差的 \(L^\infty\)-型一阶校正，都不可能把统一常数压到 \(\frac14\) 以下；并且 \(\frac14\) 在每一离散层都达到最优。

具体地：
\begin{enumerate}
  \item 在立方余链层，定理 \ref{thm:conclusion-uniform-grid-quantized-stokes-poincare-scale} 给出尺度常数 \(h/(2k)\)；取单位尺度 \(h=1\) 与 \(k=2\) 时，得到 sharp 常数 \(1/4\)；
  \item 在超立方体边流层，推论 \ref{cor:spg-hypercube-gradient-consistency-by-square-defect} 给出
  \[
  \|\omega-df\|_\infty\le \frac14\|d\omega\|_\infty,
  \]
  且 \(\frac14\) 最优；
  \item 在连续时间跳跃链层，定理 \ref{thm:spg-hypercube-near-detailed-balance-optimal} 给出近详细平衡误差同样满足 \(\varepsilon/4\) 的 sharp 控制。
\end{enumerate}
因此，\(\frac14\) 是本项目 Stokes 正规化链条的范畴不变量，而不是某一坐标选择或某一离散化习惯造成的偶然常数。
\end{theorem}
\begin{proof}
在立方余链层，定理 \ref{thm:conclusion-uniform-grid-quantized-stokes-poincare-scale} 已给出
\[
\|H_{k,h}\|_{(\infty,h)\to(\infty,h)}\le \frac{h}{2k}.
\]
当 \(h=1\) 且 \(k=2\) 时，正规化常数就是 \(1/4\)。该算子正是由连续 de Rham 同伦与 Whitney--积分互逆对诱导而来，因此光滑层与立方余链层共享同一 sharp 来源。

在超立方体边流层，推论 \ref{cor:spg-hypercube-gradient-consistency-by-square-defect} 已经精确给出
\[
\omega=df+r,
\qquad
\|r\|_\infty\le \frac14\|d\omega\|_\infty,
\]
并且常数 \(\frac14\) 不可改进。

在连续时间跳跃链层，定理 \ref{thm:spg-hypercube-near-detailed-balance-optimal} 把同一边余链 \(\omega\) 解释为对数非平衡一形式，并给出
\[
\bigl|\omega(u,v)-(f(v)-f(u))\bigr|\le \frac14\,\|d\omega\|_\infty
\]
及其最优性。因此在近详细平衡校正中，sharp 常数仍然是 \(\frac14\)。

若某一下游离散层存在统一常数 \(c<1/4\) 的正规化原理，则把正文中达到极值的方形缺陷例沿标准立方/超立方嵌入回拉，便会在超立方体边流层或连续时间跳跃链层得到同样的 \(c\)-界，从而与推论 \ref{cor:spg-hypercube-gradient-consistency-by-square-defect} 或定理 \ref{thm:spg-hypercube-near-detailed-balance-optimal} 的 sharp 性矛盾。故 \(\frac14\) 无法被统一改进。证毕。
\end{proof}

当边界积分由代数零集合在 \((\mathbb C^\ast)^2\) 中的闭链承载时，一元 Jensen 公式可视为二维 Green--斯托克斯机制在单位圆族上的纤维化表达。
在无根穿越的参数区间内，其变分给出严格的边界积分公式，从而把 Mahler 测度的导数与椭圆曲线周期直接对接。

\begin{theorem}[Lee--Yang 型零集合的 Jensen--斯托克斯变分公式]\label{thm:conclusion-leyang-jensen-stokes-variational}
令 \(t\mapsto f_t(x,y)\in\mathbb C[x^{\pm1},y^{\pm1}]\) 为 \(C^1\) 族 Laurent 多项式。
对每个 \(t\) 与每个 \(x\in S^1\)，将 \(y\mapsto f_t(x,y)\) 视为一元多项式，并假设：
\begin{enumerate}
  \item 对所有 \(x\in S^1\)，当 \(|y|=1\) 时有 \(f_t(x,y)\neq 0\)，从而可连续选择外根分支；
  \item 令 \(\gamma_t\) 为代数曲线 \(C_t:=\{(x,y)\in(\mathbb C^*)^2:f_t(x,y)=0\}\) 上由条件 \(|x|=1\) 且取全部 \(|y|>1\) 的根分支拼成的有向 \(1\)-循环（按 \(x=e^{\mathrm i\theta}\) 的 \(\theta\) 递增取向），并且该选择在参数区间内无根穿越 \(|y|=1\)。
\end{enumerate}
定义二维 Mahler 测度
\[
m(t):=\frac1{(2\pi)^2}\int_{|x|=|y|=1}\log|f_t(x,y)|\,\mathrm d\theta_x\,\mathrm d\theta_y.
\]
则 \(m\) 可微且满足变分公式
\[
m'(t)
=-\frac1{2\pi}\,\Re\int_{\gamma_t}\frac{\partial_t f_t(x,y)}{y\,\partial_y f_t(x,y)}\,\frac{\mathrm d x}{x}.
\]
若 \(C_t\) 的光滑射影完备化为椭圆曲线 \(E_t\)，则
\[
\omega_t:=\frac{1}{y\,\partial_y f_t(x,y)}\,\frac{\mathrm d x}{x}
\]
在 \(E_t\) 上延拓为亚纯 \(1\)-形式，且 \(m'(t)\) 为其在同调类 \([\gamma_t]\) 上的实周期；因此 \(m'(t)\) 满足由 Gauss--Manin 连接诱导的 Picard--Fuchs 线性微分方程（阶至多 \(2\)）。
\end{theorem}
\begin{proof}[证明要点]
对每个固定 \(x\in S^1\) 应用一元 Jensen 公式，把 \(\int_{|y|=1}\log|f_t(x,y)|\) 表示为外根模长对数之和。
在无根穿越假设下可对外根分支逐一隐式微分，并将分支求和重写为沿 \(\gamma_t\) 的线积分，得到所示公式。
当 \(E_t\) 为椭圆曲线时，上式给出 de Rham 类的周期；周期随参数变形受 Gauss--Manin 连接控制，从而满足 Picard--Fuchs 方程。证毕。
\end{proof}

最后记录一个把“周期群的圆维缺陷”与“函数方程的边界项结构”耦合的猜想性框架，其目的在于把动力学 \(\zeta\) 的谱退化与自守 \(\varepsilon\)-因子的自由度归入同一类广义斯托克斯边界机制。

\begin{conjecture}[动力学 \texorpdfstring{$\zeta$}{zeta} 与自守 \texorpdfstring{$L$}{L}-函数之间的斯托克斯边界项桥]\label{conj:conclusion-stokes-boundary-langlands-bridge}
设 \(\varphi^t\) 为紧流形上的 \(C^r\)（\(r\) 充分大）Axiom A 流，其 Markov 编码给出有限型子移位 \((\Sigma_A,\sigma)\) 与屋顶函数 \(\tau:\Sigma_A\to\mathbb R_{>0}\)。
记周期值群的 \(\mathbb Z\)-秩为 \(\rho\)：
\[
\mathcal P_\tau:=\Bigl\{\sum_{k=0}^{n-1}\tau(\sigma^k x):\ x\in\Sigma_A,\ \sigma^n x=x\Bigr\},
\qquad
\rho:=\operatorname{rank}_{\mathbb Z}\langle \mathcal P_\tau\rangle.
\]
则存在与 \((A,\tau)\) 自然关联的代数群与自守表示 \(\pi\)，使得 Ruelle \(\zeta\) 函数
\[
\zeta_{\mathrm{Ru}}(s)=\prod_{\gamma\ \mathrm{primitive}}\bigl(1-e^{-s\ell(\gamma)}\bigr)^{-1}
\]
在适当正规化后与 \(L(\pi,s)\) 的局部因子族一致；并且其函数方程中的 \(\varepsilon\)-因子可由某一流盒紧化三维链上的 Chern--Simons 型 \(3\)-形式之广义斯托克斯边界项给出，而该边界项的自由度维数恰为 \(\rho\)。
\end{conjecture}

\begin{proposition}[随机边界采样的斯托克斯证书]\label{prop:conclusion-stokes-random-sampling-certificate}
在定理 \ref{thm:conclusion-multiscale-stokes-energy-l2-reconstruction} 的设定下固定尺度 \(m\)，并设 \(\alpha\) 有界：\(\|\alpha\|_{\infty}\le B\)。
对每个 \(C\in\mathcal C_m^{k+1}\)，在 \(\partial C\) 的 \(k\)-维测度归一化概率上独立采样 \(N\) 个点 \(X_{C,1},\dots,X_{C,N}\)。
记 \(\tau_C\) 为 \(\partial C\) 上的单位定向 \(k\)-向量场，并定义 Monte Carlo 估计量
\[
\widehat{\mathbf S}_m\alpha(C)
:=\frac{|\partial C|}{h^{k+1}}\cdot \frac1N\sum_{j=1}^N \langle \alpha(X_{C,j}),\tau_C(X_{C,j})\rangle.
\]
令胞腔数 \(M_m:=|\mathcal C_m^{k+1}|\asymp h^{-d}\)。
则对任意 \(\varepsilon>0\)，有一致偏差界
\[
\PP\!\left(
\max_{C\in\mathcal C_m^{k+1}}
\bigl|\widehat{\mathbf S}_m\alpha(C)-\mathbf S_m\alpha(C)\bigr|
\ge \varepsilon
\right)
\le
2M_m\exp\!\left(
-\frac{2N\varepsilon^2}{\left(\frac{|\partial C|}{h^{k+1}}B\right)^2}
\right).
\]
特别地，若取
\[
N\ \ge\
\frac12\left(\frac{|\partial C|}{h^{k+1}}B\right)^2\cdot
\frac{\log(2M_m/\delta)}{\varepsilon^2},
\]
则以概率至少 \(1-\delta\) 同时对所有胞腔成立
\(
\bigl|\widehat{\mathbf S}_m\alpha(C)-\mathbf S_m\alpha(C)\bigr|\le \varepsilon
\)，从而由边界采样得到对 \(\mathcal E_m(\alpha)\) 的可审计统计证书，并可与定理 \ref{thm:conclusion-multiscale-stokes-energy-l2-reconstruction} 的尺度极限合并得到 \(\|d\alpha\|_{L^2}\) 的外推控制。
\end{proposition}

\begin{proof}
对固定胞腔 \(C\)，令随机变量
\[
Y_{C,j}:=\frac{|\partial C|}{h^{k+1}}\langle \alpha(X_{C,j}),\tau_C(X_{C,j})\rangle.
\]
则 \(\{Y_{C,j}\}_{j=1}^N\) 独立同分布，且由 \(\|\alpha\|_\infty\le B\)、\(\|\tau_C\|=1\) 得
\(
|Y_{C,j}|\le \frac{|\partial C|}{h^{k+1}}B
\)。
并且由采样分布的定义
\[
\EE[Y_{C,j}]
=\frac{|\partial C|}{h^{k+1}}\cdot \frac1{|\partial C|}\int_{\partial C}\alpha
=\frac1{h^{k+1}}\int_{\partial C}\alpha
=\mathbf S_m\alpha(C).
\]
由 Hoeffding 不等式得
\[
\PP\bigl(|\widehat{\mathbf S}_m\alpha(C)-\mathbf S_m\alpha(C)|\ge\varepsilon\bigr)
\le
2\exp\!\left(
-\frac{2N\varepsilon^2}{\left(\frac{|\partial C|}{h^{k+1}}B\right)^2}
\right).
\]
对全部 \(M_m\) 个胞腔作并合界即得第一式；再解出 \(N\) 得充分条件。证毕。
\end{proof}
